% NAF 5166 — lecture modernisée du volume entier.
%
% Pass: Opus 5.5 (claude-opus-5-5), 2026-09-23 — first pass, unchecked against
% the views by a human, and an interpretation rather than a record. Where this
% file and the transcriptions differ in substance, the transcriptions
% (batch-01.fr.tex to batch-04.fr.tex) are the record.
%
% Sources read: the four batch transcriptions of this volume, in full, twice —
% once for the mathematics, once for the apparatus — with their header
% comments and the coordinator's reconciliation notes. Together they cover
% views 1–80. View 81 is the back board of the binding: there is no batch 5,
% and the volume is completely transcribed. No image was consulted, and no
% printed edition was used to fill a gap, settle a word or supply a line.
%
% What the volume turned out to be. An unsorted dossier in at least four hands,
% under a nineteenth-century title « Papiers de Lagrange » and an older wrapper
% (leaf A) that claims « notes etc de Mlle Sophie Germain ». Three groups:
%
%   v. 13, 15   title leaf with the 1888 certificate; the wrapper (leaf A)
%   v. 17–41    the memoir « Eclaircissement d'une difficulté singuliere …
%               Par J. L. Lagrange », paragraphs 1–10, leaves 1–12 with 6 bis
%               (hand 1); v. 40 a cancelled first start of leaf 11
%   v. 43       a heading « Notes de M.r Dela Grange »
%   v. 45–51    short notes in hand 1: a plane and a sphere; a series headed
%               « page 623 » on slips A and B; the arc of an ellipse
%   v. 53, 55   the plane-and-sphere note and the series redone (hand 2)
%   v. 55–56    a scrap of numerical reckoning (log-cosines, degrees), pencil
%               tables on its back
%   v. 58–60    a draft letter answering a judgement, on « l'équation … d'après
%               M.r De la Grange » (hand 3, feminine agreements)
%   v. 61       another unfinished draft, on the other side of leaf 22
%   v. 62–68    a complete letter recommending « le mémoire N° 1 » to the
%               class; Chladni, Euler, Sauveur, the programme; « votre servante »
%
% Settled by the reconciliation notes and followed here: the ink series A,
% 1–26 with 6 bis is the library's count of leaves, not the writer's
% pagination; views 60 and 61 are the two sides of leaf 22. No \folio{} is
% written, the series being in ink; sections cite views.
%
% Notation translated, each with a footnote at first use: the looped d (total
% and partial alike) → d or ∂; ϖ (and the π' of views 21–39) → ϖ'; π as 180°
% → π; « V( » radical → √; ∫ with a raised 3 → triple integral over the
% body; dϖ' dθ' sin θ' → dΩ'; « &c. » → etc.; « l », « log. » → log; « sin φ² »
% → sin²φ; « b4 » on the line → exponent; R^(n) → P_n(μ); « premieres
% dimensions de y » → first order; « degré » of an equation → order; § → §.
%
% Modern names attached where the statement coincides: Newtonian potential;
% Newton's shell theorem; Legendre polynomials (Laplace's R^(n)); the
% generating function Σ(2n+1)P_n t^n; Abel summation; single-layer potential
% and its jump relation; the direct value as the mean of the two one-sided
% limits; Poisson kernel of the ball (approximate identity); Laplace's surface
% equation for spheroids (distinguished from ΔV = 0); spherical trigonometry
% (cosine and five-part formulas); Hesse normal form; binomial series;
% incomplete elliptic integral of the second kind; log-cosine tables;
% Euler–Bernoulli beam; plate equation / Germain–Lagrange equation /
% biharmonic operator; membrane equation; sum of principal curvatures (twice
% the mean curvature); simply supported and free edges; Chladni figures;
% Sauveur's nodes. Later work is footnoted as later where it is used; the
% memoir's title names Lagrange (d. 1813), and every « later than the leaf »
% is conditional on that title, stated as such.
%
% Every computation was redone; the recomputations are this reading's own.
% Departures from the leaf, each footnoted at its place: the missing minus
% signs of v. 23 and v. 27 (and the power of a in dρ/dr at r = a, v. 27); the
% « = » missing in the second line of v. 27; « 1/r » for a²/r (v. 27); the
% claim that Σ(2n+1)R^(n) « sera nécessairement nulle » and that its terms
% cancel (the series diverges; its Abel sum is 0); the third spherical-
% trigonometry formula (v. 37), whose sign is wrong; « puissances impaires de
% sin φ » (v. 37), completed; « π' et θ' deviennent π et θ » at ψ = π (v. 39);
% the component « dans le plan r et R » (v. 19), which lies in the meridian
% plane; « dx² = a cos φ dφ » (v. 51); « b4 » for b³ (v. 49); the expansion of
% log(1 + t) at t = √(1+b) > 1 (v. 55), outside its disc; « z » for y (v. 59);
% « Monsieur d'ageer » (v. 61), read « Monsieur, d'agréer ». Steps supplied:
% the implicit r ≥ a of views 25–39; the bound that makes paragraph 9 a proof;
% the reduction of a point below the reference sphere (y < 0) to the case
% y ≥ 0; the value of the illegible digit of the log table (v. 55) by
% computation, given as a computation.
%
% Attribution. The memoir names Lagrange in its title, in the text's hand;
% whose hand it is, the transcriptions do not say. The letters of v. 58–68 are
% unsigned and undated; the writer uses feminine agreements. The wrapper's
% claim that the volume holds notes by Sophie Germain is reported as the
% wrapper's. The identification of the letters' writer is given in section IV
% as an inference, with its evidence, and is adopted nowhere else. No date is
% inferred for any piece; \dating is the catalogue's, verbatim.
\documentclass[11pt,a4paper]{article}
% The preamble shared by every transcription of the mathematicians' manuscripts in Gallica.
%
% It rests on one idea: **uncertainty compounds, it does not resolve.** A
% transcription that smooths over a doubtful word has destroyed the only thing
% it was made for — a reader must be able to tell, at every point, what was on
% the page and what was supplied. The apparatus macros below are therefore the
% critical apparatus, not decoration.
%
% The same commands are recognised by `scripts/render.mjs`, which produces the
% site's reading view, and by `scripts/tei.mjs`. A macro added here must be
% added there too, or rendering fails loudly — which is the wanted behaviour.
%
% Comments here are English, like the rest of the repository. The documents
% this preamble sets are French, and that is deliberate: see the skills.

\ProvidesPackage{germain}[2026/09/15 Transcriptions of mathematicians' manuscripts in Gallica]

\usepackage{amsmath,amssymb,amsthm}
\usepackage{mathrsfs}
% Kept from the parent preamble so the renderer's diagram support stays
% available; these manuscripts draw no commutative diagrams, and nothing here uses it.
\usepackage{tikz-cd}
\usepackage[english,french]{babel}
\usepackage{fontspec}

% --- Greek ------------------------------------------------------------------
% The text face stays fontspec's default, Latin Modern, which has no Greek:
% XeTeX drops every glyph it lacks with a line in the log and nothing on the
% page, and the Greek words the manuscripts quote vanished from the PDFs. So
% the Greek and Greek Extended blocks get a XeTeX character class of their own,
% and entering or leaving a run of it switches the family to CMU Serif — the
% same Computer Modern design, with polytonic Greek — and back. Only the
% family changes: a Greek word in \textit or \textbf keeps its shape and series.
% The transitions to and from every other class carry the tokens that class
% has towards an ordinary letter, so French spacing before « ; » or inside
% « » still holds after a Greek word. No grouping, so no brace or \endgroup
% can come out unbalanced; maths is left alone, since XeTeX inserts nothing
% there. Kept identical in `exercices.sty`.
\makeatletter
\newfontfamily\germainGreekfont{cmun}[
  Extension=.otf, NFSSFamily=germaingreek,
  UprightFont=*rm, ItalicFont=*ti, SlantedFont=*sl,
  BoldFont=*bx, BoldItalicFont=*bi, BoldSlantedFont=*bl]
\newXeTeXintercharclass\germain@greek
\def\germain@greekrange#1#2{%
  \@tempcnta=#1\relax
  \loop
    \XeTeXcharclass\@tempcnta=\germain@greek
    \ifnum\@tempcnta<#2\advance\@tempcnta\@ne\repeat}
\germain@greekrange{"0370}{"03FF}
\germain@greekrange{"1F00}{"1FFF}
\def\germain@greekprev{\rmdefault}
\def\germain@greekin{%
  \edef\germain@greekprev{\f@family}\fontfamily{germaingreek}\selectfont}
\def\germain@greekout{\fontfamily{\germain@greekprev}\selectfont}
\def\germain@greekjoin{%
  \XeTeXinterchartokenstate=\@ne
  \@tempcnta=\z@
  \loop
    \ifnum\@tempcnta=\germain@greek\else
      \XeTeXinterchartoks\germain@greek\@tempcnta=\expandafter{%
        \expandafter\germain@greekout\the\XeTeXinterchartoks\z@\@tempcnta}%
      \XeTeXinterchartoks\@tempcnta\germain@greek=\expandafter{%
        \the\XeTeXinterchartoks\@tempcnta\z@\germain@greekin}%
    \fi
    \ifnum\@tempcnta<4095\advance\@tempcnta\@ne\repeat}
% babel-french sets its own classes, and saves and restores the interchar
% state, each time a language is selected: join again after it does.
\addto\extrasfrench{\germain@greekjoin}
\addto\extrasenglish{\germain@greekjoin}
\AtBeginDocument{\germain@greekjoin}
\makeatother

\usepackage{geometry}
\geometry{a4paper,margin=25mm,marginparwidth=28mm}
\usepackage{marginnote}
\usepackage{xcolor}
\usepackage[normalem]{ulem}
\setlength{\emergencystretch}{3em}
\usepackage{hyperref}
\hypersetup{colorlinks=true,linkcolor=black,urlcolor=black,pdfborder={0 0 0}}

% --- Batch metadata ---------------------------------------------------------
% Taken as they stand by the rendering script, which shows them at the head of
% the reading view. They fix what the file claims to cover. The macro names are
% the parent project's, so that the scripts stay shared; \folder here means the
% volume's slug — `fr-9115` — and \pages counts Gallica views.

\newcommand{\folder}[1]{\gdef\grothFolder{#1}}
\newcommand{\batch}[1]{\gdef\grothBatch{#1}}
\newcommand{\pages}[2]{\gdef\grothFirst{#1}\gdef\grothLast{#2}}
\newcommand{\foldertitle}[1]{\gdef\grothFoldertitle{#1}}

% The holder's own shelfmark, printed as they write it: « Français 9115 ».
\newcommand{\shelfmark}[1]{\gdef\grothShelfmark{#1}}

% The dating, copied verbatim from the holder's catalogue — never inferred
% here. « XIXe siècle » is the BnF's; a pass that dates a leaf from its content
% says so in a \note{}, as its own inference.
\newcommand{\dating}[1]{\gdef\grothDating{#1}}

% The status notice, printed in the title block and repeated as a diagonal
% watermark on every page of the PDF. The manuscripts are in the public
% domain; what this line declares is not a right but a quality — an
% unverified first pass by a machine. The skills write the line; a file
% without it gets no watermark, so leaving it out is visible in the output.
\newcommand{\watermark}[1]{\gdef\grothWatermark{#1}}

\gdef\grothFolder{?}\gdef\grothBatch{}\gdef\grothFirst{?}\gdef\grothLast{?}%
\gdef\grothFoldertitle{}\gdef\grothShelfmark{}\gdef\grothDating{}\gdef\grothWatermark{}

% --- The résumé -------------------------------------------------------------
\newenvironment{resume}
  {\small\setlength{\parskip}{0.45em}}
  {\par\medskip\hrule\medskip}

% --- Keywords ---------------------------------------------------------------
% In English, closing the modernised reading's résumé: the modern vocabulary
% under which to search for what the leaves construct. The manifest extracts
% them to tag the volume — this line is the single source of the tags.
\newcommand{\keywords}[1]{\par\smallskip\noindent
  {\footnotesize\textsc{Keywords} --- \textit{#1}}\par}

% --- The critical apparatus -------------------------------------------------

% The Gallica view number, in the margin. This is the anchor that lets a
% disagreement over a formula be settled by looking at *one* image rather than
% a volume of 750. The site uses it to turn the facsimile as the transcription
% is scrolled. A view is one scanned image: usually one side of a leaf.
\newcommand{\page}[1]{%
  \marginnote{\footnotesize\color{gray!70}\textsf{v.\,#1}}%
  \label{page:#1}\ignorespaces}

% The BnF's pencilled foliation, where the leaf carries it and a pass can read
% it: « 198r ». This is what the literature cites — « ff. 198r–208v » — and
% the one line that turns a citation into a view. Recorded, never inferred:
% a folio a pass cannot read is not written.
\newcommand{\folio}[1]{%
  \marginnote{\footnotesize\color{gray!50}\textsf{f.\,#1}}\ignorespaces}

% The modernised reading's coarser counterpart to \page{}: which views a
% section is drawn from.
\newcommand{\pagerange}[2]{%
  \marginnote{\footnotesize\color{gray!70}\textsf{v.\,#1--#2}}%
  \ignorespaces}

% Illegible: never guessed.
\newcommand{\ill}{\textcolor{red!60!black}{[\,\ldots\,]}}

% A doubtful reading: the word is given, and flagged as uncertain.
\newcommand{\uncertain}[1]{\uline{#1}}

% An editorial addition — a missing letter, an implied word.
\newcommand{\add}[1]{\textcolor{blue!50!black}{[#1]}}

% Struck out by the writer. What was crossed out is part of the document.
\newcommand{\struck}[1]{\sout{#1}}

% The transcriber's note, on the state of the page or the mathematics.
\newcommand{\note}[1]{\par\smallskip{\small\color{gray!60!black}\textit{[#1]}}\par\smallskip}

% A marginal note of hers, distinct from the body of the text.
\newcommand{\marginal}[1]{\par\smallskip\noindent\hspace{1em}%
  {\small\color{gray!60!black}#1}\par\smallskip}

% --- Title ------------------------------------------------------------------
% The title block is French, like the documents themselves.

\AddToHook{shipout/background}{%
  \ifx\grothWatermark\empty\else
    \begin{tikzpicture}[remember picture, overlay]
      \node[rotate=52, scale=3.4, align=center, text=gray!18,
            font=\sffamily\bfseries]
        at (current page.center) {\grothWatermark};
    \end{tikzpicture}%
  \fi}

\AtBeginDocument{%
  \begin{center}
    {\large\bfseries Manuscrits de mathématiciens (Gallica) ---
      \ifx\grothShelfmark\empty\grothFolder\else\grothShelfmark\fi}\\[2pt]
    {\itshape\grothFoldertitle}\\[4pt]
    {\small vues \grothFirst--\grothLast
      \ifx\grothBatch\empty\else\ (lot \grothBatch)\fi}\\[2pt]
    \ifx\grothDating\empty\else
      {\small\color{gray!60!black}Datation du catalogue : \grothDating}\\[2pt]
    \fi
    \ifx\grothWatermark\empty\else
      {\small\bfseries\color{red!45!gray}\grothWatermark}\\[2pt]
    \fi
    {\footnotesize\color{gray!60!black}Transcription établie d'après les images IIIF de Gallica.
     Source gallica.bnf.fr / Bibliothèque nationale de France.\\
     Les lectures incertaines sont soulignées, les illisibles notées [\,\ldots\,],
     les ajouts entre crochets ; \textsf{v.} numérote les vues de Gallica,
     \textsf{f.} la foliotation au crayon de la BnF.}
  \end{center}
  \vspace{1em}}

\endinput


\folder{naf-5166}
\shelfmark{NAF 5166}
\pages{1}{80}
\dating{1701-1900}
\watermark{Lecture automatique\\interprétation non vérifiée}
\foldertitle{Papiers du géomètre J.-L. LAGRANGE (1813). NAF 5166 — lecture modernisée du volume entier}

\begin{document}

\section*{Résumé}

\begin{resume}
Ce volume réunit, sous un titre du XIX\textsuperscript{e} siècle, « Papiers de
Lagrange », trois ensembles de feuillets qui n'ont pas été classés ensemble.
Le premier est un mémoire de douze feuillets, « Éclaircissement d'une
difficulté singulière qui se rencontre dans le calcul de l'attraction des
sphéroïdes très peu différents de la sphère », signé dans son titre « Par J. L.
Lagrange ». Viennent ensuite quelques notes de calcul courtes, sous une
chemise « Notes de M\textsuperscript{r} Dela Grange » : la distance d'un point
à un plan, une série de logarithmes, l'arc d'ellipse, et un brouillon de
calculs numériques. Les derniers feuillets sont des lettres et des brouillons
de lettres, sans nom ni date, d'une personne qui écrit au féminin et qui défend
son propre mémoire sur les vibrations des surfaces élastiques. La chemise la
plus ancienne annonce des « notes etc de M\textsuperscript{lle} Sophie Germain ».
C'est ce qu'affirme la chemise ; les lettres ne le disent pas.

Le mémoire part d'un calcul de Laplace. On cherche l'attraction qu'exerce un
corps homogène presque sphérique, de rayon \(a(1+y')\) avec \(y'\) petit, sur
un point de sa propre surface. Laplace en avait tiré une relation simple entre
le potentiel \(V\) et sa dérivée radiale au point attiré :
\(\tfrac12 V + a\,\partial V/\partial r = -\tfrac23\pi a^2\). L'auteur annonce
d'entrée ce qu'il cherche. Le théorème de Laplace est « très remarquable par sa
simplicité autant que par sa généralité », mais il « donne lieu à une difficulté
singulière qui paroît n'avoir encore été remarquée par personne ». Voici cette
difficulté. Si l'on dérive sous le signe d'intégration puis qu'on se place à la
surface, la fonction à intégrer est identiquement nulle, et le calcul donne
zéro. Pour une couche d'épaisseur constante, pourtant, le calcul direct donne
\(-2\pi a^2 y'\). L'auteur appelle cela « un paradoxe dans le Calcul
intégral ».

Pour le lever, il refait le calcul « sans rien négliger ». Il laisse le point
attiré à une distance \(r\) du centre plus grande que \(a\), et ne fait \(r=a\)
qu'à la fin. La fonction à intégrer vaut alors \((r^2-a^2)\) fois une fonction
qui devient infinie comme \(1/(r-a)\) près du point attiré, et le produit ne
tend pas vers zéro. L'auteur le calcule d'abord pour une couche uniforme, puis
en général par une intégration par parties. Il trouve \(-2\pi a^2 y\), et
retombe, par un chemin qu'il juge rigoureux, sur l'équation de Laplace. Il a
raison, et sa démonstration est bonne pour le point pris du côté extérieur,
pourvu qu'on y ajoute une majoration qu'il ne donne pas. Dans le langage
d'aujourd'hui, on a échangé une limite et une intégrale sans en avoir le droit.
La fonction \((r^2-a^2)\rho^3\) est, à un facteur près, le noyau de Poisson de
la boule : quand \(r\) tend vers \(a\), elle se concentre au point attiré. Au
passage, l'auteur affirme que la série \(1+3P_1(\mu)+5P_2(\mu)+\dots\), bâtie sur
les polynômes de Legendre, est nulle. C'est faux au sens propre, car elle
diverge. C'est vrai au sens d'Abel.

Les lettres se rapportent, par leur vocabulaire (« la classe », « le
programme », les « commissaires », un « mémoire N° 1 » placé sous une
épigraphe de Bacon), à un concours de l'Institut. Par leur sujet, la théorie
des plaques élastiques vibrantes comparée aux figures de Chladni, elles se
rapportent au prix des surfaces élastiques. L'une répond à un jugement.
M\textsuperscript{r} De la Grange, en suivant « mon hypothèse », a trouvé une
équation du quatrième ordre, et l'écrivante remarque qu'elle est la somme de
deux équations qu'on obtient en remplaçant, dans l'équation de la lame
élastique, la courbure \(1/r\) par la somme des courbures principales
\(1/r+1/r'\). Cette somme est \(p\,\partial_t^2 z + k^2\Delta^2 z = 0\),
l'équation des plaques vibrantes, avec l'opérateur biharmonique. Elle sait,
écrit-elle, que « ce n'est pas ainsi qu'a du operer M\textsuperscript{r} De la
Grange ». L'autre lettre plaide pour son mémoire, point par point : l'accord
avec les expériences de Chladni, le sens exact des mots « appuyé » et
« libre », la différence entre les nœuds d'une lame et ceux d'une corde. Sur
ses propres raisons, elle écrit : « je n'avois guere de confiance dans mon
travail j'y avois été entraîné par une analogie qui me sembloit frappante ».

Où cela mène, et quels noms chercher. Pour le mémoire : l'attraction des
sphéroïdes, le potentiel de simple couche et son saut à la traversée de la
couche, le noyau de Poisson, la sommation d'Abel, et la dérivation sous le
signe somme. Pour les notes : la forme normale de Hesse, la série du binôme et
l'intégrale elliptique de seconde espèce. Pour les lettres : l'équation des
plaques de Germain–Lagrange, l'opérateur biharmonique, la poutre
d'Euler–Bernoulli, les figures de Chladni et les lignes nodales. La dernière
section de cette lecture dit ce que les feuillets permettent d'affirmer sur les
auteurs, et ce qu'ils ne permettent pas.

\keywords{attraction of spheroids, Newtonian potential, Legendre polynomials, divergent series, Abel summation, single-layer potential, jump relations, Poisson kernel, interchange of limit and integral, differentiation under the integral sign, spherical trigonometry, distance from a point to a plane, Hesse normal form, binomial series, elliptic integral of the second kind, arc length of an ellipse, logarithmic trigonometric tables, vibrating plates, Germain–Lagrange plate equation, biharmonic operator, Euler–Bernoulli beam, membrane equation, sum of principal curvatures, Chladni figures, nodal lines, boundary conditions of plates}
\end{resume}

\section*{Ce que le volume contient}

\pagerange{1}{80}
Quatre-vingts vues ont été transcrites. La vue 81 est le plat inférieur de la
reliure, où rien n'est écrit. Il n'y a donc pas de cinquième lot, et le volume
est transcrit en entier. Les vues 1 et 2 sont le plat et le contreplat
supérieurs, où se trouve l'étiquette « FR. / NOUV. ACQ. / 5166 ». Les vues 3 à
12 et 70 à 79 sont des feuillets de garde modernes, blancs, et la vue 80 est le
contreplat inférieur. Plusieurs versos sont blancs eux aussi : ce sont les vues
14, 16, 18, 20, 22 à 38 paires, 42, 44, 46, 48, 50, 52, 54, 57 et 69, que les
transcriptions ont examinées sans les transcrire.

La vue 13 est un feuillet de titre calligraphié, « Papiers / de / Lagrange ».
Il porte le certificat de la bibliothèque : « Volume de 26 Feuillets / plus les
Feuillets A préliminaire \& 6 \textsuperscript{bis}. / 27 Septembre 1888. », et
la mention « Libri 1856 ». La vue 15 est le feuillet « A », une ancienne chemise.
On y lit « Lagrange. » en grande écriture, puis, d'une cursive rapide du
XIX\textsuperscript{e} siècle, « Lagrange / mémoires et notes \emph{autographes} /
avec des notes etc de / M\textsuperscript{lle} Sophie Germain ».\footnote{La
transcription lit « mémoires » avec doute (la finale est un simple trait, et
le pluriel n'est pas sûr) ; elle lit aussi « etc » avec doute (trois lettres
rapides). Sur la même chemise : « N. a. fr. 5166. », « Libri. N\textsuperscript{o}.
1856. », un « 1856 » au crayon, « R. c. 8070 (91) ». Précédé de « No. »,
« 1856 » est un numéro de la collection Libri et non une année. C'est une
inférence de la transcription, reprise ici. « R. c. 8070 (91) » et le cachet
« ACQ. 8070 (LIBRI) » de la vue 17 sont des marques d'acquisition. Rien de cela
n'est de la main des textes.} Aucune de ces mentions n'est de la main des
textes. Ce sont des notes de possesseur et de bibliothèque, et cette lecture
ne s'en sert que pour ce qu'elles disent du volume.

Une série de chiffres à l'encre court en haut à droite des rectos. Elle va de
« A. » (vue 15) à « 1 » et « 2 » (vues 17 et 19), puis de « 3 » à « 11 », avec
« 6+ » et « 6 bis » (vues 21 à 39), de « 12 » à « 22 » (vues 41 à 60) et de
« 23 » à « 26 » (vues 62 à 68). C'est la foliotation de la bibliothèque : elle
traverse quatre mains sans rupture et répond exactement au certificat de 1888.
Ce n'est pas une pagination de l'auteur, même si elle commence à 3 dans les vues
21 à 40 et y répète le 6. Comme elle est à l'encre, les transcriptions n'en
tirent aucun folio, et cette lecture cite les vues.\footnote{Le « + » qui suit
le 6 (vue 27) et une petite croix près du 7 (vue 31) restent inexpliqués. En
marge de la vue 17, une autre main a entouré d'un cercle « 13 feuillets / le
6\textsuperscript{me} double ». Ce compte s'accorde avec ce que les
transcriptions ont trouvé : le mémoire occupe les feuillets 1 à 12 et le
feuillet 6 bis, soit treize feuillets, dont le sixième est « double ». Le
rapprochement est de cette lecture.} Les autres chiffres sont d'une autre
nature. « A » et « B », à côté de 15 et de 16, apparient deux fiches d'un même
calcul. Les numéros de paragraphe du mémoire, de 1 à 10, sont cités dans le
texte sous la forme « n°. 3 ». Il y a aussi « page 623 » et « Pag. 623 », qui
renvoient à un ouvrage non nommé, et « § 3 », « § 6 », « N° 1 » dans les
lettres. Un « I 65 » encadré, d'une plume plus fine, se lit dans la marge de la
vue 37 ; son sens reste inconnu. Les vues 60 et 61 sont les deux faces d'un même
feuillet, le 22.

\begin{itemize}
  \item \textbf{Vues 17 à 41 — le mémoire} « Éclaircissement d'une difficulté
    singulière… », « Par J. L. Lagrange ». C'est un brouillon, avec des
    ratures, en paragraphes numérotés de 1 à 10, d'une cursive rapide et
    appuyée (main 1). Il occupe les rectos seulement. Au dos du feuillet 11
    (vue 40), tête-bêche, se trouve un premier départ de ce feuillet, biffé.
  \item \textbf{Vue 43 — une chemise}, « Notes de M\textsuperscript{r} Dela
    Grange », d'une écriture fine et posée.
  \item \textbf{Vues 45 à 56 — des notes de calcul} : un plan et une sphère
    (vue 45), une série en \(b\) sur deux fiches « page 623 » (vues 47 et 49),
    l'arc d'ellipse (vue 51), tous de la main 1. Les deux premiers calculs
    sont refaits d'une main soignée et arrondie (main 2, vues 53 et 55). Il y a
    aussi un brouillon de calculs numériques à l'encre, avec des tables au
    crayon au dos (vues 55 et 56).
  \item \textbf{Vues 58 à 68 — des lettres}, d'une ou de plusieurs mains
    penchées (main 3 aux vues 58 à 60 ; celle des vues 61 à 68 emploie le
    \emph{s} long). Les vues 58 à 60 sont un brouillon de réponse à un
    jugement, la vue 61 un autre brouillon, inachevé, et les vues 62 à 68 une
    lettre complète qui recommande « le mémoire N° 1 ».\footnote{Les mains 2 et
    3 ont en commun un \emph{d} à haste droite et un 5 en forme de \emph{ſ}.
    Qu'elles soient une seule main à deux vitesses, et que l'une ou l'autre
    soit celle des vues 61 à 68, les transcriptions ne le tranchent pas. La
    chemise de la vue 43 est peut-être de la main 2, sans que ce soit établi.}
\end{itemize}

\section*{I. « Éclaircissement d'une difficulté singulière » : l'attraction d'un sphéroïde en un point de sa surface (vues 17 à 41)}

\pagerange{17}{41}
\subsection*{Ce que l'auteur annonce, et les notations}

Le mémoire s'ouvre sur un historique. « D'alembert est le premier qui ait donné
le calcul de l'attraction des sphéroïdes non elliptiques mais peu differens de
la sphere », dans les deuxième et troisième volumes de ses \emph{Recherches sur
le système du monde}. Laplace a repris la question « d'une maniere nouvelle et
plus generale », dans les \emph{Mémoires de l'Académie des sciences} de 1782 et
dans le second volume de la \emph{Mécanique céleste}. Sa théorie repose sur
« un beau théoreme très remarquable par sa simplicité autant que par sa
généralité ; mais ce theoreme donne lieu à une difficulté singuliere qui paroît
n'avoir encore été remarquée par personne et qui merite d'être
examinée ».\footnote{Entre « et » et « troisieme », un mot biffé d'un trait
épais, surmonté d'un second mot biffé ; ni l'un ni l'autre ne se lit. Dans
« 1782 », le 7 est barré ; le millésime est le même à l'article 1. « paroît » :
la voyelle après \emph{par} peut être un \emph{o} ou un \emph{a}. Le titre
tient trois lignes, puis vient la ligne d'auteur, en retrait, de la même main
et de la même encre que le texte. Le cachet de la Bibliothèque nationale est
posé sur les mots « très remarquable ».}

L'article 1 fixe les notations « comme dans les Memoires de 1782 p. 134 et dans
la Mecanique celeste (tome 2) ».\footnote{Le numéro d'article est un trait
vertical à empattements, qui peut se lire « 1 » ou « I ». « (tome 2) » est écrit
dans la marge et appelé par un signe d'insertion. Un mot court biffé, après
« l'angle que », ne se lit pas.} Le point attiré \(P\) a pour coordonnées
sphériques \((r,\theta,\varpi)\) par rapport au centre du sphéroïde et à un axe
fixe. La molécule attirante \(M\) a pour coordonnées \((R,\theta',\varpi')\).
Si \(\gamma\) est l'angle des deux rayons, on pose
\[
\mu=\cos\gamma=\cos\theta\cos\theta'+\sin\theta\sin\theta'\cos(\varpi-\varpi'),
\qquad |PM|=\sqrt{r^2-2rR\mu+R^2}.
\]
L'élément de masse est \(R^2\,dR\,d\Omega'\), avec
\(d\Omega'=\sin\theta'\,d\theta'\,d\varpi'\) l'élément d'aire de la sphère
unité, et la « somme des molecules divisées par leur distances au point
attiré » est
\[
V(r,\theta,\varpi)=\iiint_{\text{corps}}\frac{R^2\,dR\,d\Omega'}{\sqrt{r^2-2rR\mu+R^2}},
\]
où \(R\) va de \(0\) à la surface, \(\varpi'\) de \(0\) à \(2\pi\) et \(\theta'\)
de \(0\) à \(\pi\).\footnote{Le feuillet écrit l'angle \(\varpi\) d'un \(\pi\)
au trait supérieur recourbé, et l'intégrale triple d'un seul signe \(\int\)
suivi d'un 3 en exposant. Le radical est un grand \emph{V} suivi d'une
parenthèse ; les angles sont en degrés (« \(360^\circ\) », « \(180^\circ\) »).
Le \(d\) bouclé de la main sert aux différentielles totales comme partielles ;
on écrit ici \(\partial\) quand la dérivée est partielle. En marge, après
\(\theta'\), « ou à \(\mu\) » est lu avec doute. Plusieurs passages sont biffés
et ne se lisent pas : une ligne commençant par « ou \(R^2dR\) », un début de
formule « \(V=\int\) », une ligne et demie après « \(180^\circ\) ; », où ne se
lisent que « et que », « depuis » et « jusqu'à », et une fraction commencée
entre la deuxième et la troisième composante. Une grosse tache couvre le signe
\(=\) de la formule de \(V\). « leur distances » : ainsi.} La densité est prise
égale à 1 sans que le feuillet le dise : le corps est supposé homogène.
\(V\) est ce qu'on appelle aujourd'hui le potentiel newtonien du
corps.\footnote{Le mot « potentiel » est plus tardif. Green l'emploie en 1828,
Gauss en 1840. Le titre du mémoire nomme Lagrange, mort en 1813 : si le titre
dit vrai, le feuillet est antérieur à ces deux textes. La « fonction \(V\) »
est celle de Laplace, que le feuillet cite.} Ses dérivées
\(\partial V/\partial r\), \(\frac1r\,\partial V/\partial\theta\) et
\(\frac{1}{r\sin\theta}\,\partial V/\partial\varpi\) sont les composantes de
l'attraction selon les directions où croissent \(r\), \(\theta\) et \(\varpi\).
L'attraction vers le centre est donc \(-\partial V/\partial r\).\footnote{Le
feuillet place la deuxième composante « perpendiculairement à ce rayon dans le
plan \(r\) et \(R\) ». En fait, \(\frac1r\,\partial V/\partial\theta\) est la
composante perpendiculaire à \(r\) dans le plan méridien, qui contient l'axe
fixe et le rayon \(r\). Le plan de \(r\) et \(R\) change avec la molécule. La
troisième composante est perpendiculaire à ce plan méridien.}

\subsection*{Le sphéroïde, une sphère et une couche}

\pagerange{21}{21}
L'article 2 prend un sphéroïde très peu différent d'une sphère de rayon \(a\) :
\(R=a(1+y')\), où \(y'\) est une fonction très petite de la direction
\((\theta',\varpi')\), dont on néglige le carré. Le potentiel est celui de la
sphère, plus celui de l'excès du sphéroïde sur la sphère. Au premier ordre,
cet excès est une couche de densité superficielle \(\sigma=ay'\) étalée sur la
sphère \(r=a\). On l'obtient en mettant \(a\) pour \(R\) et \(ay'\) pour \(dR\)
dans l'intégrale de \(V\). Avec
\[
\rho=\frac{1}{\sqrt{r^2-2ar\mu+a^2}},
\]
on a donc
\[
V=\frac{4\pi a^3}{3r}+u,\qquad u=a^3\int_{S^2}y'\,\rho\;d\Omega'.
\]
L'auteur note \(\pi\) l'angle de \(180^\circ\), et il appelle « premières
dimensions de \(y\) » ce que nous appelons le premier ordre.\footnote{Dans les
vues 21 à 39, la transcription lit \(\pi'\) la longitude de la molécule, et
\(\pi\) à la fois l'angle de \(180^\circ\) et la longitude du point attiré.
Les lots des vues 17--19 et 41, qui lisent la même main, distinguent un petit
\(\pi\) au trait recourbé (\(\varpi\), l'angle) d'un grand signe en forme de
\(\Pi\) (\(\pi\), la constante). Cette lecture écrit partout \(\varpi\),
\(\varpi'\) pour les angles et \(\pi\) pour la constante. Sur la page, « egale »
qualifie le rayon ; un mot biffé après « ainsi » (« egale » répété ?) ; un
signe ou un coefficient noyé sous une tache entre « \(u=\) » et \(a^3\), effacé
par l'auteur ; « \(ay\) » sans prime.} Le terme de la sphère,
\(4\pi a^3/3r\), n'est valable que pour \(r\ge a\). À l'intérieur de la sphère,
son potentiel est \(2\pi a^2-\frac23\pi r^2\). Le feuillet emploie la formule
extérieure sans le dire, et l'on verra que ce choix est le bon.

\subsection*{Le calcul « comme M. Laplace »}

\pagerange{23}{23}
« Je differentie maintenant comme M. Laplace par rapport à \(r\) ». En
dérivant sous le signe d'intégration, on trouve
\[
\frac{\partial V}{\partial r}=-\frac{4\pi a^3}{3r^2}+\frac{\partial u}{\partial r},
\qquad
\frac{\partial u}{\partial r}=-a^3\int_{S^2}\frac{(r-a\mu)\,y'}{(r^2-2ar\mu+a^2)^{3/2}}\,d\Omega'.
\]
Le signe moins manque sur le feuillet.\footnote{La transcription note
qu'aucun signe moins ne précède \(a^3\), et que le dénominateur est écrit avec
le radical suivi d'une parenthèse, d'un crochet et de l'exposant
\(\frac32\). Le feuillet écrit ici \(d\varpi'\,d\mu\) au lieu de
\(d\Omega'\). Ce changement suppose que le pôle des angles a été porté sur le
rayon \(r\), ce que l'article 3 fera explicitement. Le signe de l'orientation
est absorbé dans les bornes.} À la surface, \(r=a(1+y)\), où \(y\) est la
valeur de \(y'\) dans la direction du point attiré. Comme \(u\) est déjà du
premier ordre, l'auteur fait \(r=a\) dans \(u\) et dans ses dérivées. En
\(r=a\), on a \(\rho=1/\bigl(a\sqrt{2(1-\mu)}\bigr)\), et il obtient
\[
u=\frac{a^2}{\sqrt2}\int\frac{y'\,d\varpi'\,d\mu}{\sqrt{1-\mu}},\qquad
\frac{\partial u}{\partial r}\Big|_{r=a}=-\frac{a}{2\sqrt2}\int\frac{y'\,d\varpi'\,d\mu}{\sqrt{1-\mu}},
\qquad\text{donc}\qquad \frac u2+a\frac{\partial u}{\partial r}=0 .
\]
Les intégrales sont convergentes, car la singularité en \(\mu=1\) est en
\(1/\sqrt{1-\mu}\).\footnote{Sous \(u\), \(V\) et \(-a\), le dénominateur est
le petit 2 de cette main, fait comme un \emph{n} sans boucle ; le calcul
confirme la lecture 2. Le « 1 » de \((1-\mu)^{3/2}\) est empâté, et le signe
d'égalité qui suit le deuxième membre est surchargé.} Il en conclut qu'à la
surface
\[
\frac V2+a\frac{\partial V}{\partial r}=\frac{2\pi a^3}{3r}-\frac{4\pi a^4}{3r^2}
=-\frac{2\pi a^2}{3}+2\pi a^2y\qquad(r=a(1+y)),
\]
ce qui est exact pour la partie sphérique, au premier ordre. Or, dit-il,
« M. Laplace trouve \(\frac V2+a\frac{dV}{dr}=-\frac{2\pi a^2}{3}\) parce qu'il
suppose que la sphere touche le spheroide, auquel cas \(y=0\) ».\footnote{En
marge, en regard de ces équations, une note de six lignes de la main du texte
est barrée de trois grands traits croisés. On y lit : « [n']avoit pas pour
\ldots{} remplir necessairement la condition que la sphere touche la surface du
spheroide ». « avoit » et « remplir » sont douteux, et un mot n'est pas lu. Il
y a aussi, après « cette surface », un petit groupe biffé ; en fin de ligne,
après « \(y=0\). », « Cependant comme » biffé et lu avec doute, suivi d'un mot
noyé ; et, dans la marge, un mot noyé sous un trait épais.} Les deux calculs ne
diffèrent que du terme \(2\pi a^2y\). Toute la suite du mémoire tient à ce
terme.

Il faut dire tout de suite ce que vaut la dérivée calculée ainsi. Le nombre
qu'on obtient en posant \(r=a\) sous le signe \(\int\) après la dérivation est
une intégrale convergente. Ce n'est pourtant pas la dérivée de \(u\) en
\(r=a\) : en général, \(u\) n'y est pas dérivable. On verra plus bas que cette
« valeur directe » est la moyenne de deux dérivées latérales qui diffèrent.

\subsection*{Le cas d'une couche uniforme}

\pagerange{25}{25}
L'article 3 éprouve l'équation \(\frac u2+a\frac{\partial u}{\partial r}=0\)
sur le cas le plus simple, \(y'\) constant. On porte le pôle des angles sur
le rayon \(r\), « ce qui est permis, puisque cette ligne est fixe par rapport
aux memes angles ». On a alors \(\mu=\cos\theta'\), et l'on intègre en
\(\varpi'\) de \(0\) à \(2\pi\), puis en \(\mu\) de \(-1\) à \(1\) :
\[
\int_{S^2}\frac{d\Omega'}{\sqrt{r^2-2ar\mu+a^2}}
=\frac{2\pi}{ar}\Bigl[\sqrt{r^2-2ar\mu+a^2}\Bigr]_{\mu=1}^{\mu=-1}
=\frac{2\pi\bigl((r+a)-|r-a|\bigr)}{ar},
\]
qui vaut \(4\pi/r\) pour \(r\ge a\) et \(4\pi/a\) pour \(r\le a\).
Le feuillet ne donne que la première ligne. Il écrit \(r-a\) pour la racine en
\(\mu=1\), ce qui suppose \(r\ge a\) sans le dire.\footnote{« il est facile de »
termine la vue 23, au milieu de la phrase, et la vue 25 reprend. Dans la
fraction du début de la vue 25, le \(\theta'\) est récrit, comme aux vues 21
et 23 ; dans « \(\frac u2\) », le 2 est le petit 2 en forme de \emph{n}.} Il
en tire, pour \(r>a\),
\[
u=\frac{4\pi a^3y'}{r},\qquad \frac{\partial u}{\partial r}=-\frac{4\pi a^3y'}{r^2},
\qquad\text{et en }r=a:\qquad \frac u2+a\frac{\partial u}{\partial r}=-2\pi a^2y'\neq0 .
\]
C'est juste. Pour \(r<a\), en revanche, \(u=4\pi a^2y'\) est constant.
C'est le théorème de Newton : une couche sphérique uniforme n'exerce aucune
attraction à l'intérieur.\footnote{\emph{Principia}, livre I, proposition 70.
Le théorème est bien antérieur au feuillet.} La dérivée \(\partial u/\partial
r\) saute donc de \(0\) à \(-4\pi ay'\) quand on traverse la sphère \(r=a\).
Le calcul direct de l'article 2, refait ici, donne \(-2\pi ay'\), exactement la
moyenne des deux. La difficulté est déjà tout entière dans ce cas.

\subsection*{Un facteur identiquement nul, et la série de Laplace}

\pagerange{27}{29}
L'article 4 met le calcul sous une forme qui fait voir la difficulté. Comme
\(r\) n'entre que dans \(\rho\),
\[
\frac u2+a\frac{\partial u}{\partial r}
=a^3\int_{S^2}\Bigl(\frac\rho2+a\frac{\partial\rho}{\partial r}\Bigr)y'\,d\Omega',
\]
« en faisant \(r=a\) après la differentiation ». Or
\(\partial\rho/\partial r=-(r-a\mu)\rho^3\), et en \(r=a\) on trouve
\[
\rho=\frac{1}{a\sqrt{2(1-\mu)}},\qquad
\frac{\partial\rho}{\partial r}=-\frac{1}{2\sqrt2\,a^2\sqrt{1-\mu}},\qquad
\frac\rho2+a\frac{\partial\rho}{\partial r}=0\quad(\mu<1),
\]
ce que le feuillet appelle une « équation identique ».\footnote{Le feuillet
écrit \(\frac{d\rho}{dr}=\frac{r-a\mu}{(r^2-2ar\mu+a^2)^{3/2}}\), puis, en
\(r=a\), \(\rho=\frac{1}{a\sqrt2(1-\mu)}\) et
\(\frac{d\rho}{dr}=\frac{1}{2a\sqrt2(1-\mu)}\). Les deux valeurs de la dérivée
n'ont pas de signe moins, comme celle de \(du/dr\) à la vue 23, et avec elles
l'équation ne s'annulerait pas ; la transcription le relève. La barre du
radical ne couvre que le 2, alors que les valeurs demandent \(\sqrt{2(1-\mu)}\).
Enfin, telle qu'elle est transcrite, la valeur de la dérivée porte \(a\) où
il faut \(a^2\), puisque \(a\,\partial\rho/\partial r\) doit être du même ordre
que \(\rho\). Les valeurs données ici sont les valeurs exactes.} Le facteur
est nul pour \(\mu<1\) ; en \(\mu=1\), il n'est pas défini.

Le feuillet le vérifie ensuite, « comme l'a fait M. Laplace », sur le
développement de \(\rho\) en série descendante. Les \(R^{(n)}\) de Laplace
sont ici des fonctions de \(\mu\) seul. Ce sont les polynômes de Legendre
\(P_n(\mu)\) : pour \(r>a\),
\[
\begin{aligned}
\rho&=\sum_{n\ge0}\frac{a^n}{r^{n+1}}P_n(\mu),\qquad
r\frac{\partial\rho}{\partial r}=-\sum_{n\ge0}(n+1)\frac{a^n}{r^{n+1}}P_n(\mu),\\
\frac\rho2+r\frac{\partial\rho}{\partial r}&=-\frac12\sum_{n\ge0}(2n+1)\frac{a^n}{r^{n+1}}P_n(\mu).
\end{aligned}
\]
Le nom des polynômes vient après les travaux où ils apparaissent :
Legendre les introduit dans ses recherches sur l'attraction des sphéroïdes,
Laplace les note \(R^{(n)}\). Là où le feuillet écrit \(R^{(n)}\), le
développement coïncide exactement avec celui de \(P_n\).\footnote{Dans la
deuxième ligne du feuillet, un simple trait sépare
\(\frac{r\,d\rho}{dr}\) de la suite, sans signe d'égalité, et tous les termes
sont précédés de « \(-\) ». On lit
« \(\frac{r\,d\rho}{dr}-\frac{R^{(0)}}{r}-\frac{2aR^{(1)}}{r^2}-\frac{3a^2R^{(2)}}{r^3}+\&c.\) »,
qui est l'égalité donnée ici. Le 3 de « \(3aR^{(1)}\) » est empâté ; le
coefficient suivant est lu 5 avec doute (il pourrait être un 3), mais la
ligne d'après écrit nettement « \(5R^{(2)}\) ». « \&c. » est écrit d'un signe
en forme d'alpha suivi de deux points ; on écrit « etc. ».} En \(r=a\), le
feuillet écrit
\[
\frac\rho2+a\frac{\partial\rho}{\partial r}=-\frac1{2a}\bigl(R^{(0)}+3R^{(1)}+5R^{(2)}+\text{etc.}\bigr),
\]
« serie dans laquelle tous les termes doivent se detruire d'eux memes ; ce
qu'on trouve en effet après la substitution des valeurs de \(R^{(0)}\),
\(R^{(1)}\) \&c. en fonctions de \(\mu\) ».

Pour le prouver sans substituer, l'auteur remarque que
\(\rho\sqrt r=\bigl(r+\frac{a^2}{r}-2a\mu\bigr)^{-1/2}\).\footnote{Le feuillet
écrit « \(r+\frac1r\) » : le numérateur est un petit trait vertical, lu 1 ; la
ligne suivante écrit \(\frac{a^2}{r}\) au même endroit, qui est la valeur
juste.} Il dérive, multiplie par \(\sqrt r\), et obtient
\[
\frac\rho2+r\frac{\partial\rho}{\partial r}=-\frac{r^2-a^2}{2}\,\rho^3,
\]
d'où, en posant \(t=a/r<1\), l'identité
\[
\sum_{n\ge0}(2n+1)P_n(\mu)\,t^n=\frac{1-t^2}{(1-2t\mu+t^2)^{3/2}} .
\]
Cette identité est juste pour \(|t|<1\). C'est la fonction génératrice de la
suite \((2n+1)P_n\), qu'on tire de celle de Legendre,
\(\sum P_n t^n=(1-2t\mu+t^2)^{-1/2}\), en lui appliquant
\(1+2t\,\frac{d}{dt}\).\footnote{Le feuillet l'écrit sous la forme
« \(\frac{R^{(0)}}{r}+\frac{3aR^{(1)}}{r^2}+\ldots=\frac{1-a^2/r^2}{1-2a\mu/r+a^2/r^2}\times
\bigl(\frac{R^{(0)}}{r}+\frac{aR^{(1)}}{r^2}+\ldots\bigr)\) », ce qui revient au
même. Les deux dénominateurs de la première ligne de la vue 29 commencent par
le petit 2 en forme de \emph{n}. Entre le 3 et le \(R\) de « \(3R^{(1)}\) »,
un trait d'attaque en forme de parenthèse ouvrante, que rien ne ferme.} Le
feuillet conclut : « par consequent en faisant \(r=a\) la serie
\(R^{(0)}+3R^{(1)}+5R^{(2)}+\&c.\) sera necessairement nulle ».

Cette conclusion n'est pas juste telle qu'elle est écrite. La série
\(\sum(2n+1)P_n(\mu)\) diverge pour tout \(\mu\in[-1,1]\) : ses termes ne
tendent pas vers zéro, et leur module croît comme \(\sqrt n\) pour
\(-1<\mu<1\). En \(\mu=0\), par exemple, les termes non nuls valent \(1\),
\(-\tfrac52\), \(\tfrac{27}{8}\), \(-\tfrac{65}{16}\)… Les termes ne « se
détruisent » donc pas après substitution. Ce qui est vrai, c'est la limite de
la série entière :
\[
\lim_{t\to1^-}\sum_{n\ge0}(2n+1)P_n(\mu)\,t^n=0\quad\text{pour }\mu\ne1,
\qquad =+\infty\quad\text{pour }\mu=1 .
\]
La série a pour somme 0 au sens d'Abel, et elle n'a pas de somme au sens
ordinaire.\footnote{La sommation d'Abel tire son nom du mémoire d'Abel sur la
série du binôme (1826). Si le titre du mémoire dit vrai, le feuillet est
antérieur, et antérieur aussi à la distinction, due à Cauchy (1821), entre
séries convergentes et divergentes.} Faire \(r=a\) dans la série, c'est donc
déjà faire le passage à la limite que tout le mémoire va examiner.

\subsection*{Le paradoxe, et le calcul « sans rien négliger »}

\pagerange{31}{33}
L'article 5 énonce le paradoxe. La fonction \(\frac\rho2+a\frac{\partial
\rho}{\partial r}\) de \(\mu\) est toujours identiquement nulle. Son intégrale,
multipliée par \(y'\,d\Omega'\), devrait donc être nulle « par les principes du
calcul integral, suivant lesquels l'integrale est regardée comme la somme de
toutes les valeurs successives de la differentielle ». Or « Cependant [nous]
venons de voir que dans le cas le plus simple où \(y'\) est constant on a
\(\frac u2+a\frac{du}{dr}=-2\pi a^2y'\) ; ce qui est un paradoxe dans le Calcul
integral ».\footnote{Le « nous » attendu après « Cependant » n'est ni au pied
de la vue 29 ni en tête de la vue 31. Le mot biffé après « est » est court et
couvert d'un trait épais.}

L'article 6 donne l'explication : « je ferai d'abord le calcul sans rien
negliger ».\footnote{« d'abord », après « je vais », est rayé ; il revient à la
ligne suivante.} Pour \(r\ne a\), la fonction à intégrer est régulière, et la
dérivation sous le signe \(\int\) est légitime :
\[
\frac{\partial u}{\partial r}=a^3\int_{S^2}\frac{\partial\rho}{\partial r}\,y'\,d\Omega',
\qquad
r\frac{\partial\rho}{\partial r}=-r(r-a\mu)\rho^3=-\frac\rho2-\frac{(r^2-a^2)\rho^3}{2},
\]
la dernière égalité venant de \(r(r-a\mu)=\frac12(r^2-2ar\mu+a^2)+\frac12(r^2-a^2)\).
Donc
\[
\frac u2+r\frac{\partial u}{\partial r}=-\frac{a^3(r^2-a^2)}{2}\int_{S^2}\rho^3\,y'\,d\Omega' .
\]
C'est exact.\footnote{Ces trois lignes de calcul sont très corrigées. Dans
\(\frac{d\rho}{dr}\), une tache biffée sépare le signe moins de la fraction, et
le signe entre \(r^2\) et \(2ar\mu\) est surchargé ; il est lu moins avec
doute, ce qu'exige le calcul. Dans la ligne de \(\frac{r\,d\rho}{dr}\), des
taches couvrent le début du premier numérateur, un signe après le premier
terme et le début du numérateur \(r^2-a^2\), et chaque signe moins du dernier
membre est écrit sur une tache. Dans la dernière équation, un coefficient
d'abord écrit (une fraction) est couvert de hachures. Ici le 2 sous \(u\) est
bien formé, ce qui confirme la lecture des vues 23 à 29.} L'auteur voit
d'abord le terme « disparoître » en \(r=a\), puis il se reprend : puisque
\(\frac u2+a\frac{\partial u}{\partial r}\) n'est pas toujours nul, « il faut
que le terme multiplié par \(r^2-a^2\) ne disparoisse pas toujours en faisant
\(r=a\) ». Autrement dit, \(\int\rho^3y'\,d\Omega'\) doit devenir infinie comme
\(1/(r^2-a^2)\).

L'article 7 le vérifie pour \(y'\) constant. On porte le pôle sur \(r\),
comme au n° 3 (\(\mu=\cos\theta'\), \(d\mu=-\sin\theta'\,d\theta'\)). On
intègre en \(\varpi'\), puis en \(\mu\), avec pour primitive
\(\int(r^2-2ar\mu+a^2)^{-3/2}d\mu=\frac{1}{ar\sqrt{r^2-2ar\mu+a^2}}\), et
pour \(r>a\)
\[
\int_{S^2}\rho^3\,d\Omega'=\frac{2\pi}{ar}\Bigl(\frac1{r-a}-\frac1{r+a}\Bigr)=\frac{4\pi}{r(r^2-a^2)},
\qquad
\frac u2+r\frac{\partial u}{\partial r}=-\frac{2\pi a^3y'}{r}\ \xrightarrow[r\to a]{}\ -2\pi a^2y',
\]
« comme on l'a trouvée dans le n°. 3 ».\footnote{Avant \(\rho^3\), au début de
l'intégrale, une lettre est noyée sous une tache (un premier \(\rho\) ?). Devant
\(\frac{4\pi}{r(r^2-a^2)}\) et devant \(\frac{4\pi y'}{r(r^2-a^2)}\), une petite
tache carrée, peut-être un signe effacé. Le premier second membre de
l'équation finale, où l'on devine \(\pi\), \(a\) et \(y'\), est couvert de
hachures et barré. Le renvoi au n° 3 vise la valeur \(-2\pi a^2y'\) de la
vue 25.} Le calcul est juste, et il vaut « quelle que soit la valeur de
\(r\) » plus grande que \(a\), ce que le feuillet sous-entend.

\subsection*{Le cas général : le pôle transporté, et une intégration par parties}

\pagerange{35}{40}
L'article 8 traite une fonction \(y'\) quelconque des sinus et cosinus de
\(\varpi'\) et \(\theta'\). On porte encore le pôle sur la ligne \(r\), et
l'on nomme \(\varphi\) et \(\psi\) la longitude et la colatitude de la
molécule autour de ce nouveau pôle. On a
\(d\Omega'=\sin\psi\,d\psi\,d\varphi\) et \(\mu=\cos\psi\), de sorte qu'il
faut intégrer \(\rho^3y'\,d\varphi\,d\mu\) pour \(\varphi\) de \(0\) à
\(2\pi\) et \(\mu\) de \(-1\) à \(1\).\footnote{Les deux nouveaux angles sont
faits, l'un d'une boucle traversée d'une haste (\(\varphi\)), l'autre d'une
croix dont la branche droite remonte en crochet (\(\psi\)). La même forme de
croix sert aussi pour la capitale \(Y\) des vues 37 et 39 ; les deux ne se
séparent que par le sens. Il y a aussi un mot biffé avant « \(\cos\psi\) »,
« le » biffé devant « l'element », et de petites lettres serrées au-dessus de
« dans un axe », non lues. Le premier \(\theta\) de la dernière ligne de la
vue 35 porte un petit trait empâté, qui n'est pas un prime. Sous le « 9 » du
feuillet, un petit signe en forme de \emph{c}.} Mais \(y'\) est donné en
fonction des anciens angles. Il faut donc exprimer ceux-ci par les nouveaux, au
moyen du triangle sphérique qui a pour sommets le pôle fixe, le point attiré
et la molécule. Ses côtés sont \(\theta\), \(\theta'\), \(\psi\) ; l'angle au
pôle est \(\varpi-\varpi'\), opposé à \(\psi\), et l'angle au point attiré
est \(\varphi\), opposé à \(\theta'\). En mesurant \(\varphi\) à partir de
l'arc qui va du point attiré vers le pôle fixe, les formules de la
trigonométrie sphérique donnent
\[
\begin{aligned}
\cos\theta'&=\cos\theta\cos\psi+\sin\theta\sin\psi\cos\varphi,\\
\sin\theta'\sin(\varpi-\varpi')&=\sin\psi\sin\varphi,\\
\sin\theta'\cos(\varpi-\varpi')&=\sin\theta\cos\psi-\cos\theta\sin\psi\cos\varphi .
\end{aligned}
\]
La première est la formule des cosinus. La deuxième est celle des sinus, au
sens de \(\varphi\) près. La troisième est la formule dite « des cinq
éléments ».\footnote{Le feuillet écrit la troisième avec
« \(+\cos\theta\sin\psi\cos\varphi\) ». Ce signe est faux dans toute
convention compatible avec la première formule. Si l'on prend \(\varphi=0\),
la molécule est sur l'arc qui va vers le pôle, et l'on a \(\theta'=\theta-\psi\),
\(\varpi'=\varpi\) ; le premier membre vaut alors \(\sin(\theta-\psi)\), et non
\(\sin(\theta+\psi)\). L'erreur est sans effet sur la suite, qui n'utilise que
la forme polynomiale de ces expressions.} Si \(y'\) est un polynôme en
\(\cos\theta'\), \(\sin\theta'\sin\varpi'\) et \(\sin\theta'\cos\varpi'\),
c'est-à-dire en les cosinus directeurs de la molécule, il devient un polynôme
en \(\cos\psi\), \(\sin\psi\cos\varphi\) et \(\sin\psi\sin\varphi\).
L'intégration en \(\varphi\) sur un tour annule tout monôme de degré impair en
\(\cos\varphi\) ou en \(\sin\varphi\). Il ne reste donc que des puissances
paires de \(\sin\psi=\sqrt{1-\mu^2}\), et la fonction à intégrer en \(\mu\) est
un polynôme en \(\mu\) multiplié par \(\rho^3\), que l'on intègre en termes
finis.\footnote{Le feuillet dit que l'intégration en \(\varphi\) « fera
disparoitre toutes les puissances impaires de \(\sin\varphi\) ». C'est vrai,
mais incomplet : les puissances impaires de \(\cos\varphi\) disparaissent
aussi. Ce qu'il faut pour conclure, c'est qu'il ne reste que des puissances
paires de \(\sin\psi\). Un petit signe précède « \(\sin\theta'\sin\pi'\) ».
Dans la marge de droite, en regard de la ligne « Ainsi la formule… », se trouve
« I 65 » encadré, d'une plume plus fine, avec un crochet ouvrant « [ » au début
de la ligne. Aucun crochet fermant n'a été vu, et le feuillet ne dit pas ce que
marquent ces signes.} « C'est ainsi que d'Alembert a calculé l'attraction des
spheroides peu differens de la sphere ». Mais l'auteur n'a besoin que de la
réduction à \(-\rho^3y'\,d\varphi\,d\mu\).\footnote{Un trait court et épais
est biffé devant l'intégrale. Après « \(d\varphi\,d\mu\). », deux mots sont
rayés d'un trait continu ; le premier se lit « depuis » avec doute.}

L'article 9 intègre par parties en \(\mu\). Comme \(\partial\rho/\partial\mu
=ar\rho^3\), on a \(\int\rho^3d\mu=\rho/(ar)\). En \(\mu=1\), \(y'\) prend la
valeur \(y\) du point attiré, et \(\rho=1/(r-a)\) pour \(r>a\). En \(\mu=-1\),
on nomme \(Y\) la valeur de \(y'\), et \(\rho=1/(r+a)\). Ainsi, pour \(\varphi\)
fixé,
\[
\int_{-1}^{1}\rho^3y'\,d\mu=\frac{y}{ar(r-a)}-\frac{Y}{ar(r+a)}-\frac1{ar}\int_{-1}^{1}\frac{\partial y'}{\partial\mu}\,\rho\,d\mu .
\]
En \(\psi=0\) et en \(\psi=\pi\), \(y'\) ne dépend pas de \(\varphi\).
L'intégration en \(\varphi\) multiplie donc les deux termes tout intégrés par
\(2\pi\).\footnote{Le feuillet dit qu'en \(\psi=0\) et \(\psi=\pi\) « \(\pi'\)
et \(\theta'\) deviennent \(\pi\) et \(\theta\) ». C'est vrai en \(\psi=0\).
En \(\psi=\pi\), la molécule est au point antipode, où
\(\theta'=\pi-\theta\) et \(\varpi'=\varpi+\pi\). Ce qui importe, et qui est
vrai dans les deux cas, c'est que \(y'\) y est indépendant de \(\varphi\). Le
feuillet intègre de \(\mu=1\) à \(\mu=-1\), et son dernier terme porte alors le
signe \(+\) ; l'orientation choisie ici donne le signe \(-\), et le résultat
est le même.} En reportant dans l'équation du n° 6, on obtient
\[
\frac u2+r\frac{\partial u}{\partial r}
=-\frac{\pi a^2(r+a)\,y}{r}+\frac{\pi a^2(r-a)\,Y}{r}
+\frac{a^2(r^2-a^2)}{2r}\int_0^{2\pi}\!\!\int_{-1}^{1}\frac{\partial y'}{\partial\mu}\,\rho\,d\mu\,d\varphi ,
\]
et, en faisant \(r=a\),
\[
\frac u2+a\frac{\partial u}{\partial r}=-2\pi a^2y,
\]
« comme dans le cas où \(y\) est constant ».\footnote{\(Y\) est une capitale
barrée dans les deux premières fractions de la vue 39. Dans l'avant-dernière
équation, il prend la forme d'un 4 bouclé, qui est aussi celle de \(\psi\). Le
mot entre « l'equation » et « du n\textsuperscript{o}. 6 » n'a pas été lu
(« avec ce » ? « ancienne » ?). Deux mots courts sont rayés, avant et après
« n'ajouteroit ».} Le terme en \(Y\) s'annule avec \(r-a\). Le dernier terme
s'annule si l'intégrale double reste bornée quand \(r\to a\). Le feuillet
l'affirme dans sa dernière phrase, qui court de la vue 39 à la vue 41 : « Il
est bon de remarquer qu'une nouvelle integration par parties executée sur la
differentielle \(\frac{dy'}{d\mu}\rho\,d\mu\) n'ajouteroit rien à l'equation
que nous venons de trouver, car l'integrale de \(\rho\,d\mu\) etant
\(-\frac{1}{ar\rho}=-\frac{\sqrt{r^2-2ar\mu+a^2}}{ar}\) il n'en peut resulter
aucun terme où le facteur \(r-a\), qui devient nul lorsque \(r=a\), soit
detruit par la division ».\footnote{La vue 39 s'arrête sur « rien », sans
ponctuation, et la vue 41 s'ouvre sur « à l'equation que nous venons de
trouver ». La suite de la phrase, le sujet, les notations et la main font du
feuillet 12 la fin du mémoire ; c'est la conclusion du coordinateur des
transcriptions. Dans la vue 41, « \(\rho\) » est une lettre bouclée à longue
descendante, et la lecture \emph{p} n'est pas exclue. « soit » est récrit sur
un premier mot. La primitive est juste :
\(\partial_\mu\sqrt{r^2-2ar\mu+a^2}=-ar\rho\).}

La remarque est juste, et elle se complète ainsi. Pour \(r\ge a\), on a
\(r^2-2ar\mu+a^2=(r-a)^2+2ar(1-\mu)\ge 2a^2(1-\mu)\), donc
\(\rho\le1/\bigl(a\sqrt{2(1-\mu)}\bigr)\), qui est intégrable. Il faut encore
borner \(\partial y'/\partial\mu\). Pour \(\varphi\) fixé, cette dérivée peut
devenir infinie comme \(1/\sqrt{1-\mu}\) près du pôle, à cause du facteur
\(1/\sin\psi\), sauf si le gradient de \(y\) s'annule au point attiré. On
intègre donc d'abord en \(\varphi\) : la moyenne
\(\bar y(\mu)=\frac1{2\pi}\int_0^{2\pi}y'\,d\varphi\) d'un polynôme en les
cosinus directeurs, comme le suppose le feuillet, est un polynôme en \(\mu\),
et sa dérivée est bornée. L'intégrale double est alors bornée uniformément en
\(r\ge a\), et le dernier terme est \(O(r-a)\).\footnote{Même sans intégrer
d'abord en \(\varphi\), il suffit que \(y'\) soit de classe \(C^1\). Alors
\(|\partial y'/\partial\mu|\le C/\sqrt{1-\mu^2}\), l'intégrale double est
\(O\bigl(\log\frac1{r-a}\bigr)\), et le dernier terme reste
\(O\bigl((r-a)\log\frac1{r-a}\bigr)\), qui tend vers zéro. Cette majoration est
de la présente lecture.} Avec cette majoration, les articles 6 à 9 démontrent
l'égalité \(\lim_{r\to a^+}\bigl(\frac u2+r\frac{\partial u}{\partial r}\bigr)
=-2\pi a^2y\).

Au dos du feuillet 11 (vue 40), tête-bêche, trois lignes biffées de quatre
longs traits obliques sont un premier départ de ce feuillet. Elles portent un
« 11 » barré, les mots de la tête de la vue 39, et une première forme des
termes tout intégrés, \(\frac{Y(r+a)}{ar}-\frac{y(r\pm a)}{ar}\), récrite en
\(\frac{Y}{ar(r+a)}-\frac{y}{ar(r-a)}\), qui est la forme juste et celle
qu'a gardée la vue 39.\footnote{Les premiers mots de la deuxième ligne sont
noyés d'encre ; le signe du second binôme est douteux ; le dernier mot,
« ainsi », porte un trait qui lui est propre. À droite, un croquis à la
plume : deux arcs presque parallèles qui se rejoignent en pointe, coupés d'un
petit trait, sans aucune lettre. Le feuillet ne dit pas ce qu'il représente,
et cette lecture ne l'identifie pas.}

\subsection*{La conclusion : l'équation de M. Laplace}

\pagerange{41}{41}
L'article 10 rassemble les morceaux. Comme \(V=\frac{4\pi a^3}{3r}+u\)
(n° 2), en faisant \(r=a(1+y)\) et en négligeant les \(y^2\),
\[
\frac V2+a\frac{\partial V}{\partial r}
=-\frac{2\pi a^2}{3}+2\pi a^2y+\Bigl(\frac u2+a\frac{\partial u}{\partial r}\Bigr)
=-\frac{2\pi a^2}{3},
\]
« Comme M. Laplace l'a trouvé », dit la marge.\footnote{Dans la première
équation détachée de la vue 41, l'exposant de \(a\) au premier terme est un 3
suivi d'un chiffre noyé dans une tache, et le 4 du second est récrit sur un
autre chiffre. Le calcul se vérifie. Les deux notes marginales sont appelées
chacune par un signe d'insertion. Après « nul. », un petit signe barré ; un
long trait horizontal ferme le texte, et le quart inférieur de la page est
blanc.} La phrase suivante disait : « C'est l'equation de M. Laplace qui se
trouve ainsi demontrée rigoureusement et mise à l'abri de toute difficulté ».
Elle est biffée de trois traits.\footnote{« toute » est récrit et en partie
couvert d'une tache ; il est lu avec doute.} Le texte finit sur la remarque
que le terme \(2\pi a^2y\), « du à l'action de la sphere sur un point élevé au
dessus de sa surface de la quantité \(ay\), et par lequel l'equation que nous
avons trouvée (n° 2) differoit de celle de M. Laplace », se trouve
précisément détruit par l'intégrale de \(a^3\bigl(\frac\rho2+a\frac{\partial
\rho}{\partial r}\bigr)y'\,d\Omega'\), « dans laquelle le facteur
\(\frac\rho2+a\frac{d\rho}{dr}\) est toujours identiquement nul ».

L'équation obtenue est celle que Laplace place au fondement de sa théorie de
la figure des planètes : à la surface d'un sphéroïde homogène très peu
différent d'une sphère de rayon \(a\), on a, au premier ordre en \(y'\),
\[
\frac V2+a\,\frac{\partial V}{\partial r}=-\frac{2\pi a^2}{3} .
\]
Elle ne doit pas être confondue avec ce qu'on appelle aujourd'hui l'« équation
de Laplace », \(\Delta V=0\) hors des masses. Dans la lecture moderne, elle en
découle : on développe \(V\) extérieur en fonctions de Laplace, c'est-à-dire en
harmoniques sphériques, comme on le fait plus bas pour \(u\). Le feuillet
n'emploie pas ce développement, et sa démonstration n'en a pas besoin.

\subsection*{Ce qu'il faut en penser aujourd'hui}

Le calcul de l'auteur se laisse lire exactement en termes modernes. Au
premier ordre, \(u\) est le potentiel d'une couche simple de densité
\(\sigma=ay'\) portée par la sphère \(r=a\). Si l'on écrit
\(y'=\sum_nY_n\), avec \(Y_n\) la composante harmonique sphérique de degré
\(n\), alors
\[
u=4\pi a^2\sum_{n}\frac{Y_n}{2n+1}\Bigl(\frac ar\Bigr)^{n+1}\ (r\ge a),\qquad
u=4\pi a^2\sum_{n}\frac{Y_n}{2n+1}\Bigl(\frac ra\Bigr)^{n}\ (r\le a).
\]
Le potentiel \(u\) est continu à la traversée de la sphère, mais sa dérivée
radiale y saute de \(-4\pi\sigma=-4\pi ay\). C'est la relation de saut d'un
potentiel de simple couche. On en tire, en \(r=a\),
\[
\Bigl(\frac u2+a\frac{\partial u}{\partial r}\Bigr)_{\text{ext}}=-2\pi a^2y,\qquad
\Bigl(\frac u2+a\frac{\partial u}{\partial r}\Bigr)_{\text{int}}=+2\pi a^2y,
\]
et la « valeur directe », qu'on obtient en faisant \(r=a\) sous le signe
d'intégration, vaut \(0\), la moyenne des deux. Le calcul des articles 2 et 4,
signes moins rétablis, ne se trompe pas sur cette valeur. Il se trompe en la
prenant pour une dérivée.\footnote{La théorie du potentiel de simple couche,
avec ses relations de saut et sa valeur directe, s'est faite au
XIX\textsuperscript{e} siècle, de Poisson et Gauss à la fin du siècle. Si le
titre dit vrai, le feuillet est antérieur à ces travaux. Le nom de
« fonctions de Laplace » pour les \(Y_n\), comme celui d'« harmoniques
sphériques », est postérieur aux travaux de Laplace eux-mêmes.}

La phrase du paradoxe, « l'integrale est regardée comme la somme de toutes les
valeurs successives de la differentielle », n'est pas en défaut. Ce qui l'est,
c'est l'échange de la limite \(r\to a\) et de l'intégration. Pour \(r>a\), la
fonction
\[
K_r(\mu)=\frac{r(r^2-a^2)}{4\pi}\,\rho^3
\]
est positive et d'intégrale 1 sur la sphère, d'après l'article 7. Quand
\(r\to a\), elle tend vers zéro uniformément sur \(\mu\le1-\delta\), pour tout
\(\delta>0\), et toute sa masse se ramasse au point \(\mu=1\) : c'est, à la
normalisation près, le noyau de Poisson de la boule, une « approximation de
l'identité ». Pour \(y'\) seulement continu, on a donc
\(\int K_r\,y'\,d\Omega'\to y\), et
\[
\frac u2+r\frac{\partial u}{\partial r}=-\frac{2\pi a^3}{r}\int_{S^2}K_r\,y'\,d\Omega'
\ \xrightarrow[r\to a^+]{}\ -2\pi a^2y .
\]
C'est la démonstration la plus courte du résultat de l'article 9. La
fonction à intégrer tend bien vers zéro en chaque point \(\mu<1\), mais son
intégrale ne tend pas vers zéro, parce que la convergence n'est pas dominée.
Le « paradoxe » est cette différence.\footnote{Le noyau porte le nom de
Poisson, qui l'a étudié après 1813, et le critère de convergence dominée celui
de Lebesgue, au début du XX\textsuperscript{e} siècle. Si le titre dit vrai,
le feuillet précède l'un et l'autre. L'auteur n'en a pas besoin : il calcule
l'intégrale exactement.}

Reste à savoir laquelle des trois valeurs est la bonne. Pour le corps réel,
dont la densité est bornée, le gradient du potentiel est continu partout, et
il n'y a aucune ambiguïté. Le saut vient de la linéarisation : on a écrasé
la couche comprise entre les rayons \(a\) et \(a(1+y')\) sur la sphère
\(r=a\). Le point attiré est au bord extérieur de la matière, et c'est la
limite extérieure qu'il faut prendre. C'est celle que prend l'auteur aux
articles 6 à 9, où la racine en \(\mu=1\) vaut \(r-a\). Il suppose ainsi
\(r\ge a\) sans le dire, c'est-à-dire \(y\ge0\) au point attiré. Le cas
\(y<0\) se ramène à celui-là. Changer le rayon de la sphère de référence de
\(a\) en \(a(1+c)\), avec \(c\) petit, change \(\frac V2+a\frac{\partial
V}{\partial r}\) de \(ac\,\frac{\partial V}{\partial r}=-\frac{4\pi a^2c}{3}\)
et \(-\frac{2\pi a^2}{3}\) d'autant, au premier ordre. L'équation garde donc
la même forme, et l'on peut toujours prendre une sphère de référence
intérieure au corps, pour laquelle \(y'\ge0\). La sphère de rayon \(a(1+c)\)
elle-même la vérifie exactement, quel que soit le signe de \(c\) :
\(\frac{2\pi a^2(1+c)^2}{3}-\frac{4\pi a^2(1+c)}{3}=-\frac{2\pi a^2}{3}+O(c^2)\).
C'est la raison profonde de la note marginale biffée de la vue 23 : Laplace
n'avait pas besoin, pour que son équation soit vraie, que la sphère touche le
sphéroïde.\footnote{Cette réduction et cette vérification sont de la présente
lecture ; le feuillet n'aborde pas le cas \(y<0\).}

Le bilan est donc le suivant. Les articles 2, 4 et 5 ne sont pas une
démonstration : ils échangent sans droit une limite et une intégrale, et
l'article 4 conclut à la nullité d'une série divergente. Les articles 3 et 7
sont des calculs exacts pour \(r>a\). Les articles 6 à 9 démontrent le
résultat par la limite extérieure, à une majoration près, que la dernière
phrase indique et que l'on a complétée plus haut. L'article 10 en tire
l'équation de Laplace au premier ordre. L'auteur a biffé le mot
« rigoureusement ». Au terme près de cette majoration, il aurait pu le
garder.

\section*{II. « Notes de M\textsuperscript{r} Dela Grange » : quatre calculs courts (vues 43 à 56)}

\pagerange{43}{56}
Le feuillet 13 (vue 43), plus petit que les autres, ne porte qu'une ligne, d'une
plume fine et d'une écriture posée aux majuscules ornées : « Notes de
M\textsuperscript{r} Dela Grange ». C'est apparemment un titre de chemise pour
les feuillets qui suivent. Ce qu'il couvre exactement n'est pas écrit, et la
main n'est pas celle du mémoire.\footnote{« M\textsuperscript{r} » est un M
suivi d'un petit \emph{r} surélevé, souligné : « Monsieur », et non
« Mademoiselle ». « Dela Grange » est écrit en deux mots, le D et le G ornés.
Le 1 de « 13 » porte un long trait d'attaque qui le fait ressembler à un 7.}
Les quatre calculs qui suivent sont sans titre ni phrase d'introduction. Ils
sont écrits de la main du mémoire (vues 45 à 51), et deux sont refaits d'une
main soignée (vues 53 et 55).

\subsection*{La distance d'un point à un plan}

\pagerange{45}{53}
La vue 45 cherche le point d'un plan passant par l'origine,
\(Ax+By+Cz=0\), qui est le plus proche du point \((a,0,0)\). Elle écrit que
la sphère de centre \((a,0,0)\) et de rayon \(r\) touche le plan, c'est-à-dire
que la différentielle du carré de la distance,
\((x-a)\,dx+y\,dy+z\,dz\), s'annule pour tout déplacement dans le plan.
En éliminant \(dz=-(A\,dx+B\,dy)/C\), les coefficients de \(dx\) et de \(dy\)
doivent s'annuler séparément :
\[
x-a-\frac{zA}{C}=0,\qquad y-\frac{zB}{C}=0 .
\]
Reporté dans l'équation du plan, cela donne, avec \(N=A^2+B^2+C^2\),
\[
\begin{aligned}
&z=-\frac{ACa}{N},\quad y=-\frac{ABa}{N},\quad x=a-\frac{A^2a}{N},\\
&r^2=\frac{a^2A^2(A^2+B^2+C^2)}{N^2}=\frac{A^2a^2}{N},\qquad r=\frac{|A|\,a}{\sqrt N}.
\end{aligned}
\]
Tout est juste.\footnote{Le feuillet écrit \(r=\frac{Aa}{\sqrt{A^2+B^2+C^2}}\) :
\(A\) est pris positif sans que ce soit dit. « et diff. » : « et » suivi d'une
abréviation chargée d'un paraphe, lue ainsi. Dans « \(y^2\) » de la première
ligne, le \(y\) est récrit ; une petite tache suit l'exposant de \(a\) dans
« \(A^2a\) ». Les exposants 2 sont faits comme de petits accents circonflexes.
Le \(z\) est écrit à queue ; le cachet de la bibliothèque est posé sur la fin
de la formule de \(r\).} La note conclut que \(r/a\) est le sinus de l'angle du
plan avec l'axe des \(x\), égal à \(A/\sqrt N\), et que ce nombre est le cosinus
de l'angle que la normale au plan fait avec cet axe.\footnote{« des \(x\) » :
la transcription lit « dy » suivi d'un \(x\) isolé, comme l'article et la
lettre. Le mot « donc » de la dernière phrase est lu avec doute ; au-dessus,
dans la marge, un petit signe en forme d'accent circonflexe, sans texte qui
lui réponde.} En langage actuel, la distance du point \(P_0\) au plan
\(n\cdot X=0\) est \(|n\cdot P_0|/\|n\|\), et \((A,B,C)/\sqrt N\) est le
vecteur des cosinus directeurs de la normale.

La vue 53, de la main 2, refait le même calcul en phrases complètes, jusqu'à
la même équation différentielle. Elle en tire l'équation du plan sous la forme
\[
\frac{A}{\sqrt N}\,x+\frac{B}{\sqrt N}\,y+\frac{C}{\sqrt N}\,z=0,
\]
celle d'un plan passant par l'origine et perpendiculaire à la droite dont les
coordonnées sont proportionnelles à \(A\), \(B\) et \(C\). C'est la forme
normale de Hesse, à ce nom près, qui est plus tardif.\footnote{« le resultat
de ce calcul est de » : « est de » est biffé et « donne » écrit à la suite.
« proportionnelles » est coupé en fin de ligne et sa finale récrite. Le \(z\)
est écrit comme un 2 à queue, et la barre du dernier radical est doublée. La
plume fine rappelle celle du titre de la vue 43, sans que le rapprochement
puisse être établi.}

\subsection*{Une série « page 623 »}

\pagerange{47}{55}
Deux demi-feuillets, cotés « 15 A » et « 16 B » (vues 47 et 49), portent en tête
« page 623 », qui renvoie à un ouvrage non nommé. Ils développent en série la
quantité
\[
Q(b)=-\frac1{\sqrt{1+b}}+\log\frac{1+\sqrt{1+b}}{\sqrt b}
\]
pour \(b\) petit et positif. La quantité n'est écrite en entier qu'à la vue 55 ;
les vues 47 et 49 la dérivent.\footnote{« Je mets \(x\) à la place de \(b\) » :
la lettre remplacée, lue \(b\) avec doute, ressemble aussi au 6 des
dénominateurs ; la vue 49 et la vue 55 confirment \(b\). Le logarithme est
écrit « \(l\) » par la main 1 et « log. » par la main 2 ; il s'agit du
logarithme naturel, puisque \(\log2=1-\frac12+\frac13-\dots\) (vue 55).}
L'auteur dérive d'abord. Avec \(\frac{1}{1+\sqrt{1+x}}=\frac{\sqrt{1+x}-1}{x}\),
il vient
\[
Q'(x)=\frac12\Bigl(\frac{1}{(1+x)^{3/2}}+\frac{1}{\sqrt{1+x}\,(1+\sqrt{1+x})}-\frac1x\Bigr)
=-\frac{1}{2x(1+x)^{3/2}} .
\]
Il développe ensuite par la série du binôme, et intègre terme à terme :
\[
\begin{aligned}
-\frac{(1+x)^{-3/2}}{2x}&=-\frac12\Bigl(\frac1x-\frac32+\frac{3\cdot5}{2\cdot4}x-\frac{3\cdot5\cdot7}{2\cdot4\cdot6}x^2+\dots\Bigr),\\
Q(x)&=-\frac12\log x+\frac{3x}{4}-\frac{3\cdot5}{2\cdot4}\,\frac{x^2}{4}+\frac{3\cdot5\cdot7}{2\cdot4\cdot6}\,\frac{x^3}{6}-\dots+K .
\end{aligned}
\]
Le calcul se vérifie.\footnote{Les produits sont écrits avec des points,
« 3.5.7 / 2.4.6 ». Un trait oblique traverse la barre d'une fraction sans
paraître la biffer. La parenthèse ouverte avant \(1-\frac32x\), et celle de
la ligne « integrant », ne sont pas fermées. À la dernière ligne de la vue 47,
« \(K\) » est écrit au-dessus de « const. », biffé. Après la virgule vient
« const. = » suivi d'un mot court, le tout biffé ; la fin ne se lit pas.} La
vue 49 détermine la constante. « Ayant differentié et ensuite réintegré il
s'ensuit que cette integrale peut differer de la quantité proposée d'une
constante qu'il s'agit de trouver. » Pour \(b\to0\), on a
\(Q(b)=-1+\log2-\frac12\log b+o(1)\), et la série vaut \(-\frac12\log b+K+o(1)\).
Donc \(K=\log2-1\), et, pour \(0<b<1\),
\[
Q(b)=\log2-1-\frac12\log b+\frac{3b}{4}-\frac{3\cdot5}{2\cdot4}\,\frac{b^2}{4}+\frac{3\cdot5\cdot7}{2\cdot4\cdot6}\,\frac{b^3}{6}-\dots
\]
Le raisonnement est complet : deux fonctions qui ont la même dérivée sur
\(]0,1[\) diffèrent d'une constante, et la constante se lit à la
limite.\footnote{Le feuillet écrit le dernier terme avec « \(b4\) » : les
exposants sont écrits sur la ligne, « b2 », « b4 ». Le calcul de la vue 47 donne
\(b^3\), et c'est \(b^3\) qui est juste ; la transcription le relève. La série
n'est pas close par « \&c. ». « réintegré » : le mot commence par un \emph{r}
mal formé. Le 2 de « \(-1+l2\) » est récrit sur un autre chiffre. La série
converge pour \(|b|<1\) ; le logarithme demande \(b>0\).}

La fiche supérieure de la vue 55, de la main 2, porte « Pag. 623 ». Elle
reprend la même série par un autre chemin. Comme \(\log\frac{1+\sqrt{1+b}}{\sqrt b}
=\log(1+\sqrt{1+b})-\log\sqrt b\), il faut que
\[
-\frac1{\sqrt{1+b}}+\log(1+\sqrt{1+b})=\log2-1+\frac32\cdot\frac b2-\frac{3\cdot5}{2\cdot4}\cdot\frac{b^2}{4}+\dots,
\]
ce qui est la série précédente. Pour le vérifier, la fiche développe
\(\log(1+t)=t-\frac{t^2}{2}+\frac{t^3}{3}-\dots\) avec \(t=\sqrt{1+b}\). Elle
développe ensuite chaque puissance \((1+b)^{k/2}\) en \(b\), ainsi que
\(-(1+b)^{-1/2}\), et rappelle \(\log2=1-\frac12+\frac13-\dots\). Les
développements partiels sont justes.\footnote{\(\sqrt{1+b}=1+\frac b2-\frac{b^2}{8}\),
\((1+b)^{3/2}=1+\frac32b+\frac38b^2\), \((1+b)^{5/2}=1+\frac52b+\frac{15}{8}b^2\),
\(-(1+b)^{-1/2}=-1+\frac b2-\frac{1\cdot3}{2\cdot4}b^2+\frac{1\cdot3\cdot5}{2\cdot4\cdot6}b^3-\dots\).
La transcription lit « \(\frac{1.3}{4.2}\) », avec doute, en partie sous le
cachet ; la valeur est la même. Le « A » initial est une capitale ornée qu'on
prendrait pour un C. Le dénominateur 3 de la troisième ligne est récrit en
gras, et le signe moins devant \(\frac{1+2b+b^2}{4}\) est une tache repassée.}
Mais ce chemin n'est pas légitime tel quel : pour \(b>0\), \(t=\sqrt{1+b}>1\)
est hors du disque de convergence de la série de \(\log(1+t)\). Si l'on
regroupe les termes en \(b\), le coefficient de \(b\) devient
\(\frac12(1-1+1-\dots)\), qui diverge. Le coefficient vrai, \(\frac14\), n'en
est que la somme au sens d'Abel. La fiche s'arrête aux développements et ne
regroupe rien. Elle ne tire donc pas de conclusion fausse, mais elle ne
démontre rien non plus. Le premier chemin, celui des vues 47 et 49, est le
bon.

\subsection*{L'arc d'ellipse}

\pagerange{51}{51}
Le demi-feuillet « 17 » (vue 51) calcule l'élément d'arc d'une ellipse de
demi-axes \(a\) et \(b\), paramétrée par \(x=a\sin\varphi\) et
\(y=b\cos\varphi\). Avec \(c=\sqrt{a^2-b^2}\),
\[
dx^2+dy^2=(a^2\cos^2\varphi+b^2\sin^2\varphi)\,d\varphi^2=a^2\Bigl(1-\frac{c^2}{a^2}\sin^2\varphi\Bigr)d\varphi^2,
\qquad
\int\sqrt{dx^2+dy^2}=a\,E\Bigl(\frac ca,\varphi\Bigr),
\]
où \(E(k,\varphi)=\int_0^\varphi\sqrt{1-k^2\sin^2\phi}\;d\phi\) est
l'intégrale elliptique incomplète de seconde espèce, de module \(k\). C'est
juste.\footnote{Le feuillet écrit « \(dx^2=a\cos\varphi\,d\varphi\) » : le 2 est
écrit à gauche alors que le membre droit est la différentielle première ; on
lit ici \(dx=a\cos\varphi\,d\varphi\). Les carrés sont écrits
« \(\sin\varphi^2\) », « \(\cos\varphi^2\) ». L'angle est une lettre faite de deux
boucles, qu'on pourrait lire \emph{œ} ; la dernière formule la montre comme un
\(\varphi\). Le « \(E\) » porte un petit trait bouclé, sans doute un ornement
de la majuscule, et un signe semblable à « \& » suit \(E(c,\varphi)\). La
page commence sur une expression, sans phrase qui la rattache aux fiches
précédentes.} Le feuillet appelle \(c\) « l'excentricité ». C'est
aujourd'hui la demi-distance focale, et l'on réserve le mot d'excentricité à
\(e=c/a\), qui est précisément le module de l'intégrale. La notation
\(E(c,\varphi)\), module en premier, est celle qu'emploie Legendre. Cette lecture
n'en tire aucune date.

\subsection*{Un brouillon de calculs numériques}

\pagerange{55}{56}
La fiche inférieure de la vue 55 (« 20 ») est un morceau de papier déchiré,
couvert de chiffres à l'encre, sans un mot de texte. Le 4 y est fait comme un
A. On y trouve des soustractions en degrés et minutes, justes :
\(7^\circ25'-1^\circ39'=5^\circ46'\), et \(30^\circ04'-2^\circ44'=27^\circ20'\),
où le chiffre des minutes noyé sous une tache doit être 0 pour que la
soustraction tombe juste. On y trouve aussi
\(27^\circ20'-51'=26^\circ29'\), écrit sous un \(28^\circ11'\) biffé, qui était
la somme au lieu de la différence. Il y a encore la multiplication
\(346\times4\frac12=1557\) et l'addition \(7681+5740=13421\).\footnote{Les
vérifications sont de cette lecture, et c'est par elles que le chiffre noyé
est rendu, non par la vue. Dans la transcription figurent aussi une colonne
« \(60'\,0''\) / \(10'\,14\) / \(16\ 0\) », une ligne « 5 14 29 57 » et, à sa
droite, « 26 | 10 4 ||| 17 0 9 | 26 », qui ne sont pas interprétées ici. À
droite de deux calculs, un double croissant barré, qui a la forme du signe de
la lune. En bas à droite, tête-bêche, « 100 | 52 / 1001 », de lecture douteuse.}
Un petit tableau donne trois angles et trois nombres de huit chiffres :
« SC. \(5\ 6\ 56\) », « SC. \(7\ 10\ 48\) », « S.DAB \(84\ 57\ 18\) », en regard
de « \(1999658\ldots0\) », « \(99982668\) », « \(99983142\) ». Ce sont des
logarithmes tabulaires, de caractéristique 10 :
\[
\begin{aligned}
10+\log_{10}\cos5^\circ6'56''&=9{,}9982667,\\
10+\log_{10}\cos7^\circ10'48''&=9{,}9965810,\\
10+\log_{10}\sin84^\circ57'18''&=9{,}9983142 .
\end{aligned}
\]
« SC. » se lit donc « sinus du complément », c'est-à-dire cosinus, et
« S.DAB » « sinus de l'angle \(DAB\) ». Les deux traits croisés que la
transcription signale entre les angles et les nombres échangent les deux
premières lignes, comme il le faut ; le 8 final du deuxième nombre est un
arrondi de table. Le chiffre illisible du premier nombre est
un 1, par le calcul. Le « 46 » écrit sous le « 42 » final du troisième est
moins juste que 42. Enfin \(84^\circ57'18''-84^\circ55'35''=1'43''\).\footnote{Les
calculs de logarithmes sont de cette lecture, à la septième décimale ; ils
confirment les lectures « SC. » et « S.DAB » que la transcription proposait
comme une lecture du passage, non du feuillet. Le « 1 » initial de
« \(1999658\ldots0\) » ne répond à rien dans la valeur ; il n'est pas
interprété. « ST. » est suivi d'un trait long sans nombre.} Au dos (vue 56),
des tableaux de chiffres au crayon, en partie tête-bêche, portent des
colonnes de degrés et minutes (« 55.30 / 56 0 / 56.30 / 57.0 … »), des
capitales isolées (« G », « K », « E », « D ») et un croissant. S'y ajoutent un
petit groupe à l'encre écrit en travers, « \(140\mid60\) / \(7\parallel7.\) »,
et « 420 » souligné.\footnote{\(140\times3=420\), \(60\times7=420\). La
transcription ne donne que ce qui se lit sûrement : une colonne
« 54 / 20 4 / 46 / 11 4 / 48.42 », « 27 1 37 », « 437.38 », et une case
numérotée « 50 / 51 / 52 », où l'un des chiffres de « 47 30 3\ldots{} » ne se
lit pas.} Rien ne rattache ce brouillon aux calculs qui l'entourent. Le signe
de la lune et les angles font penser à un calcul d'astronomie. C'est une
conjecture, et le feuillet ne la confirme pas.

\section*{III. Des lettres sur les surfaces élastiques (vues 58 à 68)}

\pagerange{58}{68}
\subsection*{La réponse à un jugement : l'équation de M\textsuperscript{r} De la Grange}

Le feuillet 21, recto et verso (vues 58 et 59), et le recto du feuillet 22
(vue 60) portent un texte continu à la première personne, adressé à un
« vous ». Il est sans formule d'appel, sans date ni signature : c'est une
lettre, ou son brouillon, très corrigé. La personne qui écrit remercie son
correspondant de lui avoir fait « obtenir un jugement qui m'éclaire sur les
fautes que j'ai commises ». Elle explique son erreur par ses propres mots :
« je n'avois guere de confiance dans mon travail j'y avois été entraîné par une
analogie qui me sembloit frappante mais que je ne me croyois pas trop en état
d'apprecier ».\footnote{« étonné » et « entraîné » finissent par un \emph{é}
suivi d'un court trait de sortie, qui pourrait porter un second \emph{e} :
le genre ne se lit pas sur ces deux mots. Il se lit sur « egarée »,
« persuadée » (vue 58), « revenue », « frappée » (vue 60), où la finale
\emph{-ée} est nette. « connoître » et « entierement » sont biffés ; la
lecture de « connoître » sous le trait est incertaine.}

Ce qui la surprend, c'est « l'equation que vous m'indiquez d'après
M\textsuperscript{r} De la Grange ». Au premier coup d'œil, elle l'aurait crue
« formée de la somme des deux equations »
\[
p\,\frac{\partial^2z}{\partial t^2}+2k^2\Bigl(\frac{\partial^4z}{\partial x^4}+\frac{\partial^4z}{\partial x^2\partial y^2}\Bigr)=0,
\qquad
p\,\frac{\partial^2z}{\partial t^2}+2k^2\Bigl(\frac{\partial^4z}{\partial y^4}+\frac{\partial^4z}{\partial y^2\partial x^2}\Bigr)=0 .
\]
Leur somme, divisée par 2, est
\[
p\,\frac{\partial^2z}{\partial t^2}+k^2\Bigl(\frac{\partial^4z}{\partial x^4}+2\frac{\partial^4z}{\partial x^2\partial y^2}+\frac{\partial^4z}{\partial y^4}\Bigr)
=p\,\frac{\partial^2z}{\partial t^2}+k^2\,\Delta^2z=0,
\]
l'équation des vibrations transversales d'une plaque mince, où
\(\Delta^2=\Delta\Delta\) est l'opérateur biharmonique. L'histoire de
l'élasticité l'appelle équation de Germain–Lagrange.\footnote{Le nom est celui
que la littérature donne à l'équation ; il ne dit rien de la main de ces
feuillets. La théorie des plaques de Kirchhoff (1850) retient la même équation
et en fixe les conditions aux bords. Aucune des deux équations écrites ne
suffit seule : chacune privilégie une direction, et seule leur somme est
invariante par rotation des axes. Le « \(p\) » en tête des deux équations est
barré d'un trait oblique ou récrit sur une autre lettre, et il est lu avec
doute. Dans la seconde, le 2 de « \(2k^2\) » est récrit sur un signe \(+\), et
un petit « \(k^2\) » isolé se trouve dans l'interligne.} Elle ajoute : « Si je
me suis egarée dans les doubles integrales je ne dois pas m'étonner que
l'équation qu'a trouvée M\textsuperscript{r} De la Grange ne monte qu'au
quatrième degré ». Le « degré » est ici l'ordre de dérivation. La phrase laisse
entendre que l'équation de l'écrivante était d'un ordre plus élevé. Le
feuillet ne l'écrit pas.

Elle cite alors l'équation de sa « note », qu'elle dit déduite « des formules
de la mécanique analytique » :
\[
\int dz\,dy\int X\,dm+\int dx\,dz\;Y\,dm-2\int dx\,dy\int Z\,dm=2\,\mathrm{I}\,d^2z\,.\,dx\,dy .
\]
La transcription est trop incertaine pour qu'on l'interprète, et cette lecture
ne l'interprète pas.\footnote{Le second terme n'a qu'un signe d'intégration ;
le \(Y\) est repassé. Le signe entre « 2 » et « \(d^2z\) » est un trait
vertical épais, rendu \(\mathrm I\) faute de mieux. Le renvoi qui suit,
« p. 53 1\textsuperscript{re} édition », est de lecture douteuse : le 5 est
fait comme un \emph{ſ}.} Sa question est de méthode. Ne faut-il pas
différencier cette équation quatre fois, deux fois en \(x\) et deux fois en
\(y\) ? En faisant abstraction du terme dû à l'élasticité et en prenant
\(X=Y\), on obtient par ce procédé « l'equation connue de la surface tendue »,
\[
X\Bigl(\frac{\partial^2z}{\partial x^2}+\frac{\partial^2z}{\partial y^2}\Bigr)=2Z,
\]
c'est-à-dire \(X\,\Delta z=2Z\). C'est l'équation d'équilibre d'une membrane
tendue également dans les deux sens et chargée normalement, au facteur et à
la convention de signe près. Elle est, dit la lettre, « au fond la même que
celle donnée par M\textsuperscript{r} La Grange lui même p. 151 de la 1ère
édition ». Comment le même procédé deviendrait-il « vicieux lorsqu'il
s'applique au terme dependant de l'elasticité » ?\footnote{« par ce procédé
des differentiations successives » est écrit dans l'interligne, sans signe
d'insertion, et placé ici devant « obtient ». Les deux passages biffés qui
l'encadrent sont « a pour le mo[uvement] » et « pour le mouvement » ; « connue
de la surface tendue » est aussi biffé. « p. 151 » est de lecture douteuse,
le troisième chiffre étant chargé d'encre. La phrase finit « suivant les \(x\)
et les \(z\) » : la lettre anguleuse a la forme d'un \(z\), mais le sens veut
\(y\), et c'est ce qu'on lit ici. « directions » est biffé devant \(x\), qui
est souligné.} La question est légitime. Elle a sa réponse dans la mécanique
actuelle : la différentiation n'est pas en cause, c'est l'expression du terme
élastique qui l'est.

Le correspondant lui a fait remarquer que l'équation que « M\textsuperscript{r}
De la Grande a trouvée en suivant mon hypothese » contient celle de la lame
élastique, tandis que la sienne n'a pas cet avantage. En effet, si \(z\) ne
dépend pas de \(y\), \(\Delta^2z\) se réduit à
\(\partial^4z/\partial x^4\), et l'équation des plaques redonne celle de la
poutre d'Euler–Bernoulli, \(p\,\partial_t^2z+k^2\,\partial_x^4z=0\).
Elle s'était fait la même objection, et avait cru y répondre. La suite, à la
vue 60, est un brouillon trop repris pour qu'on en reconstruise la
construction. On y lit l'idée d'un « facteur qui compose l'équation propre à la
surface », qui n'est « pas necessairement egal à zéro », et d'un « autre
facteur » propre à la lame. On y lit aussi qu'elle avait cru que la formule,
« considerée avant les differentiations successives », renfermait le cas de la
lame, « et c'est ce que j'ai eu l'intention de prouver à la fin de la
note ».\footnote{« De la Grande » : ainsi, alors que la vue 58 et la ligne
précédente portent « Grange ». Après « je », un mot biffé illisible ;
« m'ettois » : ainsi ; « la même » est écrit au-dessus de « cette », biffé.
Dans l'interligne, au-dessus de « qu'en differentiant par rapport », biffé, on
lit « que le resultat de la differenti », dont la fin est biffée aussi ; la
phrase reste sans construction sûre. En tête de la vue 60 : « satisfesoit à
l'équation par la supposition A. de constante \(=0\) », où « A. » est de
lecture douteuse (peut-être un \emph{d} bouclé). « coëfficient » est biffé et
remplacé par « facteur ». La ligne biffée « tant que par conséquent il fut
necessaire, pour ce … cas que » contient un mot noyé. « n'étoit » est écrit
au-dessus de « n'est seroit pas », biffés, et les deux négations se suivent.
« sujet de » est lu avec doute. « renfermoit le » est écrit au-dessus de
« étoit également applicable au », biffé.}

La lettre se termine sur une remarque, avec une réserve. « Il est singulier
que l'on puisse obtenir les deux [équations] dont j'ai deja parlé et dont
celle-ci semble [la somme] en substituant tout bonnement dans l'équation de la
lame \(\frac1r+\frac1{r'}\) à la place de \(\frac1r\) », et en prenant
successivement \(x\) puis \(y\) pour variable de la lame.\footnote{Le passage
de la vue 60 est très repris : « est la », biffé ; « et dont celle-ci
semble » et « la somme », biffé, dans l'interligne ; « tout bonnement »
au-dessus de « dans » ; « mettant » et « pour l'une », biffés ; deux mots
courts biffés, illisibles, dont le premier paraît fait de deux lettres de
formule. Puis viennent « \(Z,x\) pour l'autre \(Z,y\) à la place de \(y\) et
\(Z\) », où \(Z\) est la grande lettre du « \(2Z\) » de la vue 59 ; celui de
« \(y\) et \(Z\) » est récrit. Le sens de cette dernière substitution n'est
pas sûr, et la paraphrase qu'on en donne ici est de cette lecture.} La
remarque est juste. Dans la lame, le terme élastique est
\(\partial_x^2(1/r)\), avec la courbure \(1/r\approx\partial_x^2z\). Si l'on
remplace \(1/r\) par la somme des courbures principales,
\(\frac1r+\frac1{r'}\approx\Delta z\), qui vaut deux fois la courbure moyenne
au premier ordre, on obtient \(\partial_x^2\Delta z=\partial_x^4z+
\partial_x^2\partial_y^2z\), puis \(\partial_y^2\Delta z\) de même, et leur
somme est \(\Delta^2z\). L'opérateur des plaques est le laplacien de la somme
des courbures. « Je sais bien que ce n'est pas ainsi qu'a du operer
M\textsuperscript{r} De la Grange mais cette idée m'a tellement frappée à la
vue de son équation que je n'ai pû resister à la tentation de vous la
communiquer en protestant au reste de ma parfaite soumission aux oracles
\ldots{} jugement d'un autre grand géomètre. »\footnote{« par La Grange », sans
« M\textsuperscript{r} », deux lignes plus haut. « des maîtres et » est biffé
après « oracles », et « jugement d'un autre grand » est écrit à la suite,
d'une encre plus pâle et d'un trait plus fin. La phrase qui en résulte est
incomplète. Le texte finit sur « géomètre. », suivi d'un point ; le reste de
la page est blanc, sans signature.}

\subsection*{Un second brouillon, au dos du feuillet 22}

\pagerange{61}{61}
L'autre face du feuillet 22 (vue 61) porte un brouillon de lettre, inachevé
et très corrigé. « Vous m'engagez à reprendre une tache que j'ai déja fort mal
remplie ; je m'y sens peu disposée et il faut convenir que puisque j'ai pû me
tromper si grossierement je ne suis que trop fondée à me méfier de mes forces,
cependant j'aime assez la verité pour etre encore contente de la connoitre lors
même qu'elle ne m'est pas favorable ». Le reste remercie le correspondant de
son intérêt et du « soin que vous prenez de me dissimuler le ridicule attaché
au mauvais succès ». La lettre s'arrête sur « que vous me témoignez », biffé,
sans point ni formule.\footnote{« tache » est la « tâche ». « avez la bonté
de » est biffé devant « m'engagez », dont la fin est surchargée (un \emph{z} ou
un \emph{r} corrigé). Au-dessus de « reprendre », un grand \emph{D} bouclé isolé.
« nouvelle » et « d'ailleurs » sont dans l'interligne, et « bien
sincerement », biffé dans la ligne, est récrit dans la marge. « auquel
j'attache le plus » et « d'un » sont biffés : la phrase lit « un dédomagement
grand prix », sans correction complète. « à la médiocrité » est biffé. Deux
signes brefs biffés, illisibles, suivent un « que ». « à tant de » est écrit
au-dessus de « à la bontés », et « tant » est peut-être barré.} On y lit :
« je vous prie, Monsieur \emph{d'ageer}, l'expression de ma
reconnoissance ». La transcription prend « d'ageer » pour un nom, qu'elle
n'identifie pas. La construction demande un verbe (« je vous prie … d'agréer
l'expression … et de croire que »), et la lettre des vues 62 à 68 finit par la
même formule, « en vous priant, Monsieur, d'agréer l'assurance ». Cette
lecture lit donc « Monsieur, d'agréer » et non un nom de
destinataire.\footnote{C'est un écart par rapport à la transcription, fondé
sur la construction et non sur la vue. La transcription donne les lettres
\emph{d}, apostrophe ou accent, \emph{a}, \emph{g}, \emph{e}, \emph{e} ou
\emph{c}, \emph{r}. Il manquerait un \emph{r} à « agréer ». L'image seule en
décidera.} Les transcriptions ne disent pas lequel des deux textes du feuillet
22 fut écrit le premier, ni s'ils forment un seul brouillon.

\subsection*{La lettre qui recommande le « mémoire N° 1 »}

\pagerange{62}{68}
Les feuillets 23 à 26 (vues 62 à 68) portent une lettre complète et continue,
du premier mot à la souscription, mais sans nom, sans date, sans lieu et sans
nom de destinataire. Elle est écrite d'une cursive penchée qui emploie le
\emph{s} long, et ses accords sont au féminin. Elle commence ainsi : « Je
recommande à votre protection le mémoire N\textsuperscript{o} 1 portant cette
épigraphe » :
\begin{quote}
« Sed longe maximum progressibus scientiarum, et novis pensis ac provinciis in
isdem suscipiendis obstaculum deprehenditur in desperatione hominum et
suppositione impossibilis. » \emph{Baco novum org.}
\end{quote}
L'épigraphe est tirée du \emph{Novum Organum} de Bacon.\footnote{« Mais de
loin le plus grand obstacle au progrès des sciences, et à l'entreprise de
tâches et de domaines nouveaux en elles, se trouve dans le désespoir des
hommes et dans la supposition de l'impossible. » (traduction de cette
lecture). « isdem » : ainsi, pour \emph{iisdem}. Chaque ligne de l'épigraphe
est ouverte par des guillemets en marge, et « Baco novum org. » est souligné.
« N\textsuperscript{o} 1 » : un \emph{N} à \emph{o} suscrit, puis un chiffre
bref, lu 1.} L'écrivante aurait voulu consulter son correspondant avant de
l'adopter, car elle a « un air de venterie » qui ne lui convient guère.
Elle ne voit pas « de plus forte objection contre ma théorie que
l'improbabilité d'avoir rencontré juste ».\footnote{« venterie » est lu avec
doute ; c'est la « vanterie ». « j'en avois » est écrit au-dessus de
« j'eusse », biffé, et le « en » est peut-être barré aussi.} Elle craint
« l'influence de l'opinion qu'avoit manifestée M\textsuperscript{r} De la
Grange ». La question n'a sans doute été abandonnée, écrit-elle, « que
parceque ce grand géomètre l'avoit jugée difficile », et le même préjugé
pourrait la faire condamner sans examen. Elle demande que « les commissaires
prennent la peine d'entendre mon long et ennuieux travail », bien que « le sujet
du mémoire ne me laisse pas l'espérance de vous avoir pour
juge ».\footnote{« De la Grange » est écrit en deux mots, coupés en fin de
ligne. « sentence » : le \emph{n} du milieu est repassé. « d'avantage » :
ainsi.}

Le plaidoyer suit le mémoire point par point.
\begin{itemize}
  \item \textbf{L'accord avec l'expérience.} Elle a comparé sa théorie aux
    expériences de M\textsuperscript{r} Chladni. Il reste une différence, peu
    considérable, qui n'apparaît que dans certains cas. Ce sont les cas qui
    appartiennent à une intégrale particulière « qui donne lieu à une
    équation de la forme de celle aux cordes vibrantes », quand on compare
    les sons de ses figures à ceux d'une autre intégrale particulière. Les
    sons d'une même intégrale, comparés entre eux, donnent des rapports
    « conformes à l'expérience ». Elle juge sa théorie « plus avancée, sous le
    rapport de la comparaison à l'expérience que celle des surface tendues
    sur laquelle personne n'élève de doute ».\footnote{« cette » (« de cette
    théorie ») est surchargé, peut-être biffé, et gardé. « renfermées » est
    écrit « renfermés ». « absolument » est biffé deux fois, et « mêmes » et
    « entièrement » sont biffés aussi. « d'une », « que je crois » et
    « pense » sont écrits au-dessus de mots biffés. Avant « appuyée », un mot
    biffé illisible, suivi d'un groupe bref qui ressemble à « dûe ». « des
    surface tendues » : ainsi.} Les figures sont celles que Chladni rendait
    visibles avec du sable sur les plaques vibrantes : les figures de Chladni,
    qui dessinent les lignes nodales.
  \item \textbf{Le § 3}, sur « la valeur analytique des mots \emph{appuyé} et
    \emph{libre} », employés « d'une manière tellement vague qu'il en est
    résulté, même dans le mémoire d'Euler, une grande confusion d'idée ». Ce
    sont les conditions aux bords : un bord simplement appuyé, où la flèche
    et le moment de flexion sont nuls, et un bord libre, où le moment et
    l'effort tranchant sont nuls.\footnote{« principal » (« l'objet principal
    de mes recherches ») est lu avec doute. « le » est barré devant « § 3 »,
    et « dans lequel j'ai fait examen » est traversé d'un trait, les deux
    derniers mots lus avec doute. La phrase qui reste cite le titre d'un
    paragraphe. Sur la page, les \emph{s} finaux de « mouvements », « lames »,
    « plaques » et « les » sont barrés, le pluriel ramené au singulier.}
  \item \textbf{Les lignes de repos.} Les points et lignes de repos de la lame
    et de la plaque élastiques diffèrent de ceux de la corde et de la surface
    tendue. Cette différence « tient à l'ordre des équations différentielles
    de ces corps sonores ». Elle résulte, pour la lame, de la théorie
    d'Euler, du travail de Ricati, et de l'expérience. Avec des lames assez
    longues, l'expérience montre entre les lignes de repos des intervalles
    « sensiblement inégaux », conformes à la théorie d'Euler, ce que Chladni,
    qui s'est « sans doute » servi de lames trop courtes, n'a pas
    mentionné.\footnote{« Ricati » : ainsi. La page se termine sur « enfin
    de », et la phrase continue au verso. Un mot biffé, « un » ou « en »,
    précède « servi ».} Chez la corde, le nœud de Sauveur est « une véritable
    limite analytique » : si l'on coupait la corde en son nœud, chaque moitié
    vibrerait comme avant. Chez la lame, dans « le plus grand nombre » des
    cas, les parties séparées par des points de repos ne pourraient, isolées,
    ni se ployer aux mêmes courbures ni rendre les mêmes sons. C'est juste,
    et c'est bien affaire d'ordre. Pour l'équation de la corde, qui est du
    second ordre, un nœud (\(z=0\)) est exactement la condition d'extrémité
    fixe. Pour l'équation de la lame, qui est du quatrième ordre, le moment
    \(z''\) et l'effort tranchant \(z'''\) ne s'annulent pas en général en un
    nœud. La partie isolée ne satisferait donc ni aux conditions d'un bord
    appuyé ni à celles d'un bord libre.
  \item \textbf{Le programme.} Elle demande qu'en tout état de cause la
    classe change, si la question est « encore remise », le passage de son
    programme sur les découvertes de Chladni, qui aurait rendu perceptibles
    « des \emph{nappes vibrantes} analogues aux \emph{ondes} des cordes de
    Sauveur, et des courbes \emph{d'équilibres} ou de repos auxquels
    correspondent les \emph{nœuds} ou \emph{points de repos} des mêmes
    cordes ». L'analogie qu'il suggère est précisément celle qui est fausse
    pour les lames.\footnote{« Je crois donc » est écrit au-dessus de « Il
    résulte de mes remarques », biffé. « dans le cas qu'elle soit encore
    remise » est écrit dans la marge, sur trois lignes, avec « le » biffé et,
    au-dessus de « cas », deux lettres brèves (« au » ou « ou ») de place
    incertaine. « nappes » est lu avec doute. La citation est ouverte par des
    guillemets en tête de chaque ligne, et le guillemet fermant est restitué.
    « il s'agit des expériences de » est biffé dans la parenthèse. Les mots
    en italique sont soulignés sur la page. « d'équilibres » : ainsi.}
  \item \textbf{Le § 6}, « un véritable hors d'œuvre au sujet du mémoire ».
    Il expose « un fait singulier » : en appliquant un corps mou aux côtés
    d'une plaque rectangulaire élastique, on peut tirer d'une portion aussi
    petite qu'on veut le même son que d'une portion de longueur donnée. Elle
    en donne la raison : « les termes qui doivent être nuls aux limites, sont
    l'un et l'autre multipliés par un facteur dépendant de la rigidité
    naturelle ». Au bord libre, en effet, le moment et l'effort tranchant sont
    proportionnels à la rigidité de flexion. Là où une matière non élastique
    remplace la plaque, ces deux conditions sont satisfaites d'elles-mêmes.
    Le paragraphe du mémoire n'est pas dans ce volume, et cette lecture ne
    peut pas vérifier le fait lui-même.\footnote{« hors d'œuvre » est lu avec
    doute, et le « 6 » de « § 6 » est net. « continueroient » (au feuillet
    précédent) a des lettres repassées. Après « les parties », un mot biffé
    illisible, peut-être un premier « séparées ». Sous « multipliés », un petit
    chevron sans appel.}
  \item \textbf{Le dernier chapitre} explique les différences entre la
    théorie de l'anneau d'Euler et les expériences de Chladni. « Ainsi voila
    plusieurs points parfaitement indépendans de la théorie des surfaces et
    qui sont appuyés tant sur la théorie d'Euler que sur des expériences
    faciles à répéter. »\footnote{« plusieurs » est écrit au-dessus de
    « trois », biffé, et « l'explication » après « une », biffé.}
\end{itemize}

La lettre finit sur un regret : « Combien je regrette l'avis de
M\textsuperscript{r} De la Grange ! s'il m'eut condamnée il auroit au moins
remarqué ce qui a indépendamment de la théorie principale, peut mériter
l'attention, s'il m'eut aprouvé j'aurais été sûre d'avoir raison ».
L'idée que la question est difficile, qui empêchera peut-être « que l'on ne se
donne la peine d'examiner un mémoire condamné d'avance », n'aurait alors pu lui
donner aucune crainte. Elle réclame donc la protection de son correspondant.
Puis vient la formule finale : « en vous priant, Monsieur, d'agréer l'assurance
des sentimens de d'admiration et de respect de votre servante ».\footnote{« mes
idées » est écrit au-dessus de « mon travail », biffé. Avant « Combien », deux
mots biffés, illisibles. « aurais » : la voyelle se lit plutôt \emph{a}, alors
que la main écrit ailleurs « auroit ». « mémoire » et « me donner » sont écrits
au-dessus de « travail » et de « m'inspirer », biffés. Il n'y a pas de point
après « honoré ». « respect et » est biffé devant « d'admiration » ; « et de »,
sous le cachet, est peut-être rayé. La lettre s'arrête sur « votre servante »,
sans nom.} Le regret porte sur un avis que l'écrivante n'a pas eu. La lettre
ne dit pas pourquoi elle ne l'a pas eu, et cette lecture ne le devine pas.

\section*{IV. Ce que les feuillets permettent d'affirmer}

\pagerange{13}{68}
\textbf{Le mémoire des vues 17 à 41.} Son titre porte « Par J. L. Lagrange »,
de la même main et de la même encre que le texte. C'est un brouillon, et
l'auteur y parle à la première personne. Les transcriptions ne disent pas si
la main est celle de Lagrange ; la chemise de la vue 15 parle de papiers
« autographes ». Le mémoire cite le tome 2 de la \emph{Mécanique céleste} et
lui est donc postérieur. C'est une borne que le contenu impose, pas une date,
et cette lecture n'en tire pas d'année.

\textbf{Les notes des vues 43 à 56.} La chemise dit « Notes de
M\textsuperscript{r} Dela Grange ». Les calculs des vues 45 à 51 sont de la
main du mémoire, et deux d'entre eux sont recopiés par une autre main, aux
vues 53 et 55. Le brouillon numérique est d'une autre main encore, qui
n'écrit que des chiffres. La chemise ne dit pas quels feuillets elle couvre.

\textbf{Les lettres des vues 58 à 68.} Voici ce que les feuillets montrent.
L'écrivante emploie des accords féminins et signe « votre servante ». Elle est
l'auteur d'une théorie des surfaces élastiques vibrantes et d'un mémoire
présenté à « la classe », sous le numéro 1 et une épigraphe de Bacon. Ce
mémoire compare la théorie aux expériences de Chladni, discute Euler, et
repose sur une « hypothèse » dont M\textsuperscript{r} De la Grange a tiré une
équation du quatrième ordre. Voici ce que la chemise affirme : que le volume
contient des « notes etc de M\textsuperscript{lle} Sophie Germain ». Le
vocabulaire de ces lettres (« la classe », « le programme », les
« commissaires », un mémoire numéroté sous une épigraphe) est celui d'un
concours de l'Institut. Leurs sujets, les vibrations des surfaces élastiques
comparées aux figures de Chladni, sont ceux du prix que la Première classe de
l'Institut mit au concours sur la théorie des surfaces élastiques.

L'identification de l'écrivante avec Sophie Germain est une inférence de
cette lecture, non un fait des feuillets. Plusieurs indices la soutiennent. La
chemise l'annonce. L'écrivante est une femme qui a concouru à un prix de
l'Institut sur les surfaces élastiques. Son hypothèse porte sur la somme des
courbures principales \(\frac1r+\frac1{r'}\). Lagrange a tiré de cette
hypothèse l'équation des plaques, ce qui répond à ce que l'histoire du prix
rapporte de Sophie Germain et de Lagrange.\footnote{Par exemple L. L.
Bucciarelli et N. Dworsky, \emph{Sophie Germain: An Essay in the History of
the Theory of Elasticity}, 1980. Cette lecture n'a pas consulté ce livre pour
lire les feuillets, et ne s'en sert que pour dire où l'inférence cherche sa
confirmation.} D'autres éléments manquent ou vont en sens contraire. Aucun
feuillet ne porte de nom, de date ni de signature. La chemise est d'une main du
XIX\textsuperscript{e} siècle, qui n'est pas celle des lettres. Les
transcriptions n'ont comparé ces mains à aucun autographe authentifié. Elles ne
décident pas non plus si les mains des vues 53--55, 58--60 et 61--68 sont une
seule main. L'inférence est donc forte, mais elle n'est pas établie. Ni le
destinataire, un « Monsieur » qui n'est pas nommé, ni l'autre « grand
géomètre » ne sont identifiés ici.

\textbf{Quand.} Aucun feuillet ne porte de date. Le champ de datation est
celui du catalogue, 1701--1900, et cette lecture n'en tire aucune année, ni
pour le mémoire, ni pour les notes, ni pour les lettres. Les « avant » et
« après » des notes de bas de page qui situent le mémoire par rapport à des
travaux ultérieurs sont conditionnels : ils supposent que le titre dit vrai
en nommant Lagrange.

\end{document}
